\documentclass[ecta,nameyear,draft]{econsocart}
\usepackage{amsmath,amsfonts,amssymb,amsthm,booktabs,array,longtable}
\usepackage[T1]{fontenc}
\usepackage{mathpazo,microtype,needspace,xurl}
\RequirePackage[colorlinks,citecolor=blue,linkcolor=blue,urlcolor=blue]{hyperref}
\startlocaldefs
\newtheorem{theorem}{Theorem}
\newtheorem{proposition}[theorem]{Proposition}
\newtheorem{lemma}[theorem]{Lemma}
\newtheorem{corollary}[theorem]{Corollary}
\newcommand{\E}{\mathbb E}
\newcommand{\R}{\mathbb R}
\endlocaldefs
\setlength{\emergencystretch}{3em}
\makeatletter\def\copyright@text{Prepared for referee review}\makeatother

\begin{document}
\begin{frontmatter}
\title{Neural Bellman Operators}
\runtitle{Neural Bellman Operators}
\begin{aug}
\author[id=au1,addressref={add1}]{\fnms{Qian}~\snm{QI}\ead[label=e1]{qiqian@pku.edu.cn}}
\address[id=add1]{Peking University}
\end{aug}
\begin{abstract}
This paper develops Neural Bellman Operators for policy evaluation and feasible improvement in controlled economies. A constructive continuation retains feasible action witnesses, and its native and affine--ReLU realizations preserve the same policy. Signed own-policy evaluation bands yield a centered, acquisition-aware improvement rule that retains the incumbent unless a feasible change is certified nonworsening. We prove a finite-sweep accuracy bound for the original Bellman optimum in terms of policy-evaluation, action-search and comparison errors; the evaluation-width coefficient is sharp. In the exact case, optimality follows within the finite horizon without intermediate cost increases. A source-bound investment study executes an analytically integrated final-date specialization and separately records construction, certification, actual policy costs and verification work. Both witness and conventional policies receive the same improvement. The evidence identifies own-incumbent cost reductions while generally retaining FVI's lower cost. Simultaneous cost intervals support resource and replacement decisions. The original theory, economic applications and adverse comparisons are preserved in complete development editions.
\end{abstract}
\begin{keyword}
\kwd{Dynamic programming}\kwd{Neural approximation}\kwd{Policy certification}
\kwd{Information acquisition}\kwd{Computational economics}
\end{keyword}
\end{frontmatter}
\section{Introduction}\label{sec:introduction49}
A continuation value is useful to an economist because it changes a decision. Its prediction error is not the loss of the policy actually implemented. Nor is a smaller upper bound on that loss evidence that one policy has lower cost than another. When observation, action selection and verification are costly, the numerical procedure and the economic decision have to be evaluated together.

We develop Neural Bellman Operators as a policy-evaluation and feasible-improvement method. The economic object remains a controlled economy with policy-specific continuations. The operator constructs each date's continuation using its own fitted future, returns a feasible policy, and accounts for its loss against the Bellman optimum. A general signed policy sandwich permits arbitrary fitted candidates. A constructive affine--ReLU backend additionally produces the candidate from primitive moduli, without assuming that a nonconvex training algorithm eventually succeeds. Exact compilation retains the original action witness through off-grid evaluation and state-dependent feasibility repair.

The principal question in this revision is how that construction produces a certified economic answer with finite resources. State coverage, action search and integration have different error coefficients. Tying them to one grid size conceals the allocation problem. We derive a separated primitive budget and an optimal continuous allocation for its leading work proxy, including lower resource bounds and dyadic implementation. The empirical comparison then records the complete positive-tolerance frontier of a fixed common resource ladder, not only an isolated favorable target. All failed earlier attempts belong to the work required to return a policy.

The next question is what the returned policy does. We evaluate the actual policies at the declared first crossings, even when two generators cross at different resolutions. Independent interval simulation preserves the original continuous initial and innovation laws and encloses discontinuous actors and robust capacity repair. Conditioning on construction information permits simultaneous cost inference after this selection. The result does not infer a policy ranking from two certificates and does not extrapolate three deterministic timing repetitions into an optimizer-success probability.

Two economic decisions make this distinction operational. First, a controller replacement is worthwhile only if its actual gain exceeds the replacement fee. We identify the directional gain intervals and the full break-even fee account, retaining but no longer relying on a single generous charge. Second, observation precision is chosen on the same finite bit grid on which the implemented policies are evaluated. A simultaneous upper-net-cost selector has a regret bound relative to the best precision in that catalogue. The bound holds for every nonnegative information price on one coverage event. This is a result about actual net policy cost, different from minimizing a conservative all-state loss account.

These advances do not depend on calling a classical envelope neural. A native min-plus envelope and its exactly equivalent ReLU circuit are one method with two encodings. The original-owner compiler is available to both. We compare that method with uniform multilinear fitted-value iteration (FVI) and two separately recorded conventional refinement rules. The original curvature-bisection block remained uniform; the later graded block was nonuniform. A new residual-driven comparator below records its realized geometry and all pilot work. Each generates its own future labels and receives the same feasible action menu, continuous-law integration and one-sided economic criterion. The evidence must determine the cost of the returned policy as well as certificate sharpness. No uniform neural dominance is built into the design.

\paragraph{Related methods.}
Lipschitz extension supplies the cone construction \citep{mcshane1934}, and separable distance transforms supply its linear-work nodal engine \citep{felzenszwalb2012}. Approximation in min-plus spaces has an established dynamic-programming literature: \citet{lakshminarayanan2014} study min-plus projected Bellman equations and supremum-norm errors. Bellman-inequality methods likewise construct lower value bounds and policies; \citet{wang2015} use iterated inequalities with linear or semidefinite optimization. Our argument does not claim priority for these principles. Its subject is their connection to an own-future feasible policy, original-owner continuous compilation, acquired-state implementation, complete resource allocation and direct economic comparison of the returned objects. The empirical-Bernstein component uses the bounded-sample inequality of \citet{maurer2009}; conditional selection and numerical path enclosures specify the estimand to which it is applied.

The broader NBO program includes controlled diffusions, recursive utility, endogenous preferences, temporal selves and games. Those results and their assumptions are retained in the complete development edition and its full proof supplement, together with the original training, matrix-factor and precision experiments. The present article supplies a self-contained constructive theorem--evidence chain. Its direct-cost statistical theorems concern expected discounted costs, not an undeclared extension to every nonlinear recursive application. This organization preserves the paper's economic subject while making the current contribution and the actual comparison burden explicit.

The additional improvement result addresses the economically important case where the certified construction returns a more costly policy than a conventional competitor. We preserve that incumbent and permit only actions whose actual reference advantage is certified nonpositive after the numerical and observation allowances. At the final date of the original investment model, exact continuous-law integration makes this gate executable without assuming a successful policy-value fit. Both generators receive the same repair. The resulting experiment isolates removable final-action loss from the policy gap that remains; it does not rename the compiler as a uniquely neural advantage.

\section{The economic operator and its policy account}\label{sec:core49}
Consider dates $t=0,\ldots,T$ and an invariant compact state domain $X\subset\R^d$. At date $t$, the action belongs to a nonempty compact set $A_t(x)$. The transition is $F_t(x,a,Z_{t+1})$, the current cost is $c_t(x,a)$, and the terminal cost is $g(x)$. Write $B_0=1$ and $B_t=\prod_{j<t}\beta_j$, with $\beta_j>0$. An implemented policy includes its observation rule, action selection and feasible numerical realization, not just its stored parameters.

For a continuation $v$, define
\begin{equation}\label{eq:operator49}
 Q_t(v;x,a)=c_t(x,a)+\beta_t\rho_t(v(F_t(x,a,Z))),\qquad
 \mathcal T_tv(x)=\inf_{a\in A_t(x)}Q_t(v;x,a).
\end{equation}
Throughout the deterministic comparison, the evaluations $\rho_t$ are monotone and cash invariant on a common domain containing all functions and constant translates used below. They are therefore nonexpansive in the supremum norm. Continuous objectives, effective compact feasible correspondences and measurable minimizers, or measurable arbitrarily near-minimizers with their error charged, are assumed. The stipulated conditional recursion gives Bellman comparison against all admissible adapted controls. In the experiments $\rho_t$ is conditional expectation, costs are bounded and innovations have their original continuous uniform law. The direct-cost inference in Section~\ref{sec:decisions49} is specific to expectation; it is not automatically an inference theorem for nonlinear recursive utility or games.

A Neural Bellman Operator is a construction and evaluation procedure for \eqref{eq:operator49}. It forms targets using its \emph{own} fitted $f_{t+1}$, returns a neural continuation and a feasible policy, and records an economic error allowance. It does not receive the optimal continuation as a training label. Arbitrary trained networks can be evaluated by the first result below. The subsequent constructive backend specifies how a network and its actor are obtained without assuming successful nonconvex optimization.

\subsection{Action-dependent errors and complete-policy loss}
Let $H(a)$ be a true current cost and $\widehat H(a)$ its continuation approximation. If $\widehat a$ is $\delta$-optimal for $\widehat H$ on the same feasible action set, then
\begin{equation}\label{eq:center49}
 H(\widehat a)-\inf_aH(a)\le\delta+
 \operatorname{osc}_{a}\{\widehat H(a)-H(a)\}.
\end{equation}
Indeed, compare $\widehat a$ with an arbitrarily near-optimal true action and cancel the approximate objectives. The infimum limit gives the result. Additive continuation levels cancel; action-dependent discrepancies do not. Under a change in future policies, the continuation must be transported to the new economic problem before this inequality licenses reuse. The full development edition retains the original quantitative transport, recursive-domain and economic-application results. No transport bound is presumed valid after an unaccounted change of constraints or preferences.

\begin{proposition}[One-sided policy sandwich]\label{prop:sandwich49}
Suppose, uniformly in state,
\[
 f_t-\mathcal T_tf_{t+1}\le u_t,\qquad
 \mathcal T_t^\pi f_{t+1}-f_t\le v_t,
 \qquad \ell_T\le f_T-g\le u_T.
\]
Then the actual feasible policy satisfies
\begin{equation}\label{eq:sandwich49}
 0\le J_0^\pi-V_0^*\le
 \sum_{t<T}B_t(u_t+v_t)+B_T(u_T-\ell_T).
\end{equation}
The bound is invariant to date-specific constant shifts of the fitted continuations, provided the associated residual accounts are shifted consistently.
\end{proposition}
\begin{proof}
Backward comparison and cash invariance give
\[
 V_t^*\ge f_t-\sum_{k=t}^{T-1}B_{t,k}u_k-B_{t,T}u_T,
 \quad
 J_t^\pi\le f_t+\sum_{k=t}^{T-1}B_{t,k}v_k-B_{t,T}\ell_T,
\]
where $B_{t,k}=\prod_{j=t}^{k-1}\beta_j$. The terminal claim is the stipulated enclosure; applying the Bellman or fixed-policy operator proves each inductive step. Subtraction gives \eqref{eq:sandwich49}. A shift $f_t\mapsto f_t+c_t^0$ adds $c_t^0-\beta_tc_{t+1}^0$ to $u_t$ and subtracts it from $v_t$. The terminal width is unchanged.
\end{proof}

A sharper valid residual account for the same continuation and the same deployed actor sharpens knowledge of the policy. It does not retrain that policy or change its realized cost. Conversely, subtracting two upper bounds in \eqref{eq:sandwich49} does not identify the difference in actual costs. These distinctions govern both the historical comparisons and the new experiments.

\section{Constructing and compiling a feasible neural policy}\label{sec:construction49}
Assume the primitive objective has, for feasible state--action pairs, state modulus $A_t^Q=C_{x,t}+\beta_tF_{x,t}L_{t+1}$ and action modulus $D_t=C_{a,t}+\beta_tF_{a,t}L_{t+1}$. Let
\[
 d_H(A_t(x),A_t(z))\le K_t^A\|x-z\|_1,
 \qquad L_t=A_t^Q+D_tK_t^A,\qquad L_T=L_g.
\]
The Hausdorff condition and the primitive moduli imply that $\mathcal T_tf_{t+1}$ is $L_t$-Lipschitz whenever $f_{t+1}$ is $L_{t+1}$-Lipschitz. All moduli and numerical integration procedures are supplied by the economic primitives, not fitted from successful outcomes.

At state nodes $x_{t,i}$ with covering radius $h_t$, form effective feasible $k_t$-nets. Evaluate their own-future objectives with certified error at most $e_t$. Minimize the computed labels and retain the corresponding \emph{original} action $a_{t,i}$. A Borel repair $\Gamma_{t,i}(x)$ is feasible at $x$ and satisfies
\begin{equation}\label{eq:repair49}
 \|\Gamma_{t,i}(x)-a_{t,i}\|_1\le K_t^A\|x-x_{t,i}\|_1+\zeta_t.
\end{equation}
The allowance $\zeta_t$ concerns displacement; feasibility itself may not fail. Define
\begin{equation}\label{eq:cone49}
 f_t(x)=\min_i\{y_{t,i}+L_t\|x-x_{t,i}\|_1\},\quad
 i_t(x)=\min\arg\min_i\{y_{t,i}+L_t\|x-x_{t,i}\|_1\}.
\end{equation}
Return $\pi_t(x)=\Gamma_{t,i_t(x)}(x)$. Terminal labels approximate $g$ with error $e_T$ and use the same cone construction with modulus $L_g$.

\begin{theorem}[Own-future construction with feasible witnesses]\label{thm:construct49}
The functions in \eqref{eq:cone49} have exact affine--ReLU realizations and Lipschitz moduli $L_t$ independent of label error. The returned policy is feasible and measurable. It satisfies
\begin{align}\label{eq:enclosures49}
 -e_t\le f_t-\mathcal T_tf_{t+1}&\le D_tk_t+e_t+2L_th_t,\\
 \mathcal T_t^\pi f_{t+1}-f_t&\le e_t+D_t\zeta_t,\notag
\end{align}
and hence
\begin{equation}\label{eq:constructionloss49}
 J_0^\pi-V_0^*\le
 \sum_{t<T}B_t(D_tk_t+2e_t+2L_th_t+D_t\zeta_t)
 +B_T(2e_T+2L_gh_T).
\end{equation}
With effective covers, repairs and certified queries at the prescribed accuracy, allocating each nonnegative term a share of a positive economic tolerance gives a terminating finite construction. A fixed numerical precision or a capped finite experimental ladder need not contain that allocation.
\end{theorem}
\begin{proof}
The selected node label lies between the node Bellman infimum minus $e_t$ and that infimum plus $D_tk_t+e_t$. The Lipschitz property of the infimum implies that every cone lies above it minus $e_t$. At a nearest node its cone exceeds the infimum by at most $D_tk_t+e_t+2L_th_t$. At the attaining original index, the two feasible pairs and \eqref{eq:repair49} give
\[
 Q_t(f_{t+1};x,\Gamma_{t,i}(x))
 \le y_{t,i}+e_t+L_t\|x-x_{t,i}\|_1+D_t\zeta_t.
\]
This proves \eqref{eq:enclosures49}. The terminal residual belongs to $[-e_T,e_T+2L_gh_T]$; Proposition~\ref{prop:sandwich49} gives the loss bound. A finite minimum of equally Lipschitz cones is equally Lipschitz. The identities $|z|=\operatorname{ReLU}(z)+\operatorname{ReLU}(-z)$ and $\min(u,v)=u-\operatorname{ReLU}(u-v)$ give the exact network. Finite lexicographic comparison and Borel repair give measurability. Effective refinement of the finitely many positively weighted terms yields the termination assertion.
\end{proof}

The Lipschitz-extension principle is classical \citep{mcshane1934}. A native min-plus envelope and its exact ReLU circuit, using the same labels, witnesses and ties, are \emph{one method with two encodings}. They return the same function and policy. Neither Theorem~\ref{thm:construct49} nor a faster representation proves a uniquely neural capability or a training-success probability for Adam or L-BFGS. The mathematical increment is the complete feasible, own-future construction and implementation account; the empirical question is its work and policy value relative to conventional generators.

\subsection{Original-owner compilation on a tensor cover}
On a tensor cover $\mathcal X=\prod_{j=1}^d\{x_{j,0},\ldots,x_{j,n_j}\}$, put $S=\prod_j(n_j+1)$. For arbitrary labels $y_i$, including inconsistent inexact labels, compute at each grid node $v$
\[
 (z_v,o_v)=\min_{\rm lex,i}(y_i+L\|v-x_i\|_1,i).
\]
The index $o_v$ is the original label owner, not the transformed node $v$.

\begin{theorem}[Continuous original-owner identity]\label{thm:compile49}
For a closed tensor cell $C$ with corners $\mathcal V(C)$, and $x\in C$,
\begin{align}\label{eq:compile49}
 f(x)&=\min_{v\in\mathcal V(C)}(z_v+L\|x-v\|_1),\\
 (f(x),i_*(x))&=\min_{\rm lex,v\in\mathcal V(C)}
 (y_{o_v}+L\|x-x_{o_v}\|_1,o_v).\notag
\end{align}
Coordinatewise forward and backward relaxations compute $(z,o)$ in $O(dS)$ additions and comparisons and $O(S)$ storage. After locating the cell, a query uses $O(d2^d)$ arithmetic operations. Exact rational comparisons preserve all ties, including $L=0$. Using the original owner's repaired action preserves the uncompiled policy, its cost and its certificate exactly.
\end{theorem}
\begin{proof}
The triangle inequality bounds $f(x)$ above by each corner expression. Choose the lexicographically first attaining site $i_*(x)$. In every coordinate choose the cell endpoint between $x_j$ and $x_{i_*,j}$. At this corner distances add exactly. An owner of strictly smaller transformed value would improve the objective at $x$; a tied smaller owner would contradict the original tie rule. Thus this corner has owner $i_*(x)$, proving both identities. Along a coordinate line, forward and backward relaxations respectively minimize over sites on either side, always comparing value--original-index pairs. Iterating over coordinates gives the separable distance transform and its stated work. Rational arithmetic and unchanged repair complete the claim.
\end{proof}

The nodal algorithm is the classical distance transform \citep{felzenszwalb2012}; preserving the original action through continuous off-grid queries is the policy obligation. If there are $J_{t,i}$ node actions and $M_t$ integration representatives, write $Q_t=M_t\sum_iJ_{t,i}$. Literal dense cone evaluation costs $O(\sum_t dS_{t+1}Q_t)$ in representation work. Compilation replaces this by $O(\sum_t[dS_{t+1}+Q_td2^d])$, plus cell location, primitive evaluation, integration, feasible repair, certificate and output. Rational bit lengths and state boxes crossing many cells incur additional costs. There is no dimension-free complexity assertion.

\subsection{Acquired states and conventional actors}
A $b$-bit coordinate sensor reports a dyadic cell and its midpoint $\widehat x$. On the unit cube, $\|x-\widehat x\|_1\le r(b)=d2^{-b-1}$. For the investment capacity $\kappa(x)$ below, clip the original action to the smallest capacity throughout that cell and round down to an action quantum $\Delta_t$. This is feasible at every state in the cell. Relative to an exact true-state repair, its displacement is at most $2K_t^Ar(b)+\Delta_t$. Transporting the witness cone from $\widehat x$ to $x$ adds at most $2L_tr(b)$, plus any certified selector objective error $\nu_t$. Thus the implemented-policy account adds
\begin{equation}\label{eq:acquire49}
 \sum_{t<T}B_t\{(2L_t+2D_tK_t^A)r(b)+D_t\Delta_t+\nu_t\}.
\end{equation}
No continuity of the discrete owner selector is assumed. At uncertain numerical comparisons the simulator retains all possible actors, rather than choosing a favorable one.

For conventional FVI the continuation is multilinear on its realized tensor grid. Its actual modulus is bounded by the largest coordinate edge slope. A nearest-node actor uses a feasible node minimizer and true-state repair. Let $h_t$ be the largest half-cell $\ell^1$ radius, $L_t^{\rm gr}$ the Bellman graph modulus, and $e_t$ the node query allowance. Its optimal residual upper bound is $D_tk_t+e_t+L_t^{\rm gr}h_t$. Its selected-policy upper bound is $e_t+E_t$, where
\[
 E_t=\sup_x\{y_{i(x)}+L_t^{\rm gr}\|x-x_{i(x)}\|_1-f_t(x)\}.
\]
Partitioning cells by nearest-node boundaries makes the expression separately affine in each coordinate. Its maximum is therefore at a partition vertex; the implementation encloses all such vertices. Both uniform and curvature-refined FVI receive this sharper one-sided account, not an artificially doubled baseline. Their acquired-state allowance is derived separately in \eqref{eq:fvisensor49}.

\section{Cost-directed feasible improvement}\label{sec:improvement50}
A policy certificate answers a safety question relative to the optimum. It need not answer the comparison with an installed controller. We now make actual cost reduction a constraint on the improvement step. The continuation representation, including a neural realization, supplies proposals; a centered, policy-specific evaluation determines whether to implement them. The incumbent is always an admissible fallback.

Fix a feasible implemented policy $\pi$ and bounded functions $h_0,\ldots,h_T$ in the common evaluation domain of Section~\ref{sec:core49}. Its observation and rounding rules are part of $\pi$. Suppose verified constants satisfy
\begin{equation}\label{eq:evalbands50}
 \ell_t\le \mathcal T_t^\pi h_{t+1}-h_t\le u_t\quad(t<T),
 \qquad \ell_T\le g-h_T\le u_T.
\end{equation}
Define $a_T=\ell_T$, $b_T=u_T$ and
\[
 a_t=\ell_t+\beta_t a_{t+1},\qquad
 b_t=u_t+\beta_t b_{t+1},\qquad w_t=b_t-a_t.
\]
All inequalities in this section are uniform over the true state domain. Bounds at sampled points alone do not meet this requirement.

\begin{proposition}[Centered reference advantage]\label{prop:advantage50}
Under \eqref{eq:evalbands50}, $a_t\le J_t^\pi-h_t\le b_t$. For any feasible action $v$ at state $x$, put
\[
 d_t(x,v)=Q_t(h_{t+1};x,v)-Q_t(h_{t+1};x,\pi_t(x)).
\]
If $\overline d_t(x,v)$ is a valid upper enclosure of this difference, then
\begin{equation}\label{eq:advantage50}
 Q_t(J_{t+1}^\pi;x,v)-J_t^\pi(x)
 \le \overline d_t(x,v)+\beta_t w_{t+1}.
\end{equation}
The width allowance and the action contrast are invariant to date-specific constant shifts of the fitted continuations.
\end{proposition}
\begin{proof}
Monotonicity and cash invariance propagate the two error endpoints backward. Applying them once more to each of the two feasible actions bounds each continuation-evaluation error between $\beta_ta_{t+1}$ and $\beta_tb_{t+1}$. Their difference is at most $\beta_tw_{t+1}$. Date-specific shifts translate both endpoints by the same amount and cancel in the action contrast.\end{proof}

The proposition is deliberately a reference-policy result. Substituting a bound on $h-V^*$ without controlling $h-J^\pi$ changes the statement. Likewise, accepting an action merely because it lowers $Q_t(h_{t+1})$ ignores the reference-evaluation allowance.

\begin{theorem}[Incumbent-preserving Bellman improvement]\label{thm:improve50}
Suppose a measurable feasible proposal $v_t$ and a nonnegative measurable margin $\kappa_t$ are given. At a state where a verified enclosure satisfies
\begin{equation}\label{eq:gate50}
 \overline d_t(x,v_t(x))+\beta_tw_{t+1}\le-\kappa_t(x),
\end{equation}
implement $v_t(x)$; otherwise implement the original $\pi_t(x)$. Call the resulting policy $\mu$. For a cell-based observation rule, require feasibility and \eqref{eq:gate50} for every true state in the reported cell, using the same proposed action on that cell. Then $J_t^\mu(x)\le J_t^\pi(x)$ at every date and state.

Under conditional expectation, let $D_t$ be the accepted-state set and let the state trajectory follow $\mu$. Then the stronger quantitative assertion is
\begin{equation}\label{eq:occupancy50}
 J_0^\mu(x)-J_0^\pi(x)
 \le-\E_x^\mu\sum_{t<T}B_t\kappa_t(X_t)\mathbf1_{D_t}(X_t).
\end{equation}
The qualitative assertion holds for the stipulated monotone cash-invariant conditional recursion. The occupancy formula is asserted only for expectation.
\end{theorem}
\begin{proof}
At accepted states Proposition~\ref{prop:advantage50} gives $\mathcal T_t^\mu J_{t+1}^\pi\le J_t^\pi-\kappa_t$. Elsewhere the two actions coincide, so the advantage is exactly zero, not $\beta_tw_{t+1}$. Backward monotone comparison, starting from a common terminal cost, gives the first assertion. Under expectation subtract $J_t^\pi(X_t)$ along the $\mu$ trajectory and telescope the conditional expectations. The resulting identity is the expected discounted sum of true reference advantages, each bounded above by the accepted negative margin or by zero.\end{proof}

The improvement principle is related to approximate policy iteration with safeguarded updates, including \citet{metelli2021}. We do not claim that monotone policy improvement itself is new. Here the relevant connection is between a signed own-policy evaluation account, action-centered numerical errors, original-owner construction, and the actual information-constrained action. The test is available to neural and conventional proposals alike. An exact native encoding of a ReLU continuation does not become a second independent method.

\subsection{A fully executed final-date specialization}\label{sec:terminal50}
A useful test of this result leaves every earlier decision fixed and improves only the final action. Set $h_T=g$. The continuation width is then zero, so the gate concerns the actual conditional final cost, without a fitted policy-evaluation error. For the investment economy in Section~\ref{sec:study49}, write
\[
 y_j=\frac1{16}+\frac{x_j}2+\frac{x_{j+1}}8+
       \frac{x_j(1-x_{j+1})}{16}+\gamma_ja,
 \quad s_j=(-1)^{j+1},\quad \sigma_z^2=\frac1{3072},
\]
where the cyclic convention and $\gamma_j$ are as in that section. Put $b=1/2-2\sum_jy_j/d$, $r=|\sum_js_j|/(16d)$, and
\begin{equation}\label{eq:psi50}
 \Psi_2(b,r)=
 \begin{cases}
 b_+^2,&r=0,\\
 \big((b+r)_+^3-(b-r)_+^3\big)/(6r),&r>0.
 \end{cases}
\end{equation}
Exact integration under the original continuous innovation law gives
\begin{align}\label{eq:exactterminal50}
 \E g(y+sZ)={}&\frac4d\sum_j(y_j-5/8)^2
  +\frac1{4d}\sum_j(y_j-y_{j+1})^2+2\Psi_2(b,r)\notag\\
 &+4\sigma_z^2+\frac{\sigma_z^2}{4d}\sum_j(s_j-s_{j+1})^2.
\end{align}
Thus no terminal innovation mesh is needed for this improvement step. Its computation is nevertheless charged. The formula includes the odd-dimensional shortage uncertainty; setting it to zero in dimension three would change the model.

\begin{corollary}[Acquired-cell final improvement]\label{cor:terminal50}
In the declared investment economies, final conditional cost
$H(x,a)=c(x,a)+\beta\E g(F(x,a,Z))$ is strictly convex in $a$, with second derivative at least $2p$. Fix an observation cell $C$ and the implemented incumbent action $a_C$. Choose a lattice minimizer $a_C^*$ at the reported midpoint on the cell-wide feasible set. Accept it only if an outward enclosure of $H(x,a_C^*)-H(x,a_C)$ is nonpositive for all $x\in C$. Keep the incumbent on every other cell and at all earlier dates. The resulting acquired policy weakly improves actual expected discounted cost for every initial state. This guarantee does not compare it to a different full-information incumbent.
\end{corollary}
The lattice is $\{k/4096:k\in\mathbb N_0\}$, truncated by the minimum capacity on the eight-bit cell. Convex derivative bisection identifies the two neighboring lattice actions; their actual cost contrast selects the midpoint proposal. The gate compares its cost with the incumbent on the entire cell, not just at the midpoint. Exact rational point comparison resolves a numerically ambiguous proposal tie. An uncertain true-state box may intersect several cells; the simulator retains their action hulls rather than selecting a favorable cell.

For each generator, old and repaired policies share their earlier trajectory. Their discounted cost difference therefore equals the conditional final contrast. Centering cancels state-only costs and tightens its interval enclosure. When two generators are compared, their trajectories differ and are simulated separately using common innovations. Both comparisons concern actual policies; neither is obtained by subtracting two policy-loss certificates.

\section{From Safe Improvement to Finite-Sweep Accuracy}\label{sec:finite-sweeps51}
The preceding improvement result compares an implemented policy with its own
incumbent. We now connect that comparison to the original Bellman optimum.
This distinction matters economically: a policy can improve strictly while
remaining more costly than a conventional alternative. The result below
preserves the original objective and gives explicit sufficient conditions for
closing its optimality gap; it does not infer those conditions from a small
empirical residual.

Consider the finite-horizon expectation model of the paper, with bounded
costs, feasible measurable Markov policies, and discount factors
$\beta_t\geq0$, $0\leq t<T$. Write
\[
 Q_t(h;x,a)=c_t(x,a)+\beta_t\mathbb E[h(X_{t+1})\mid x,a],\qquad
 (T_th)(x)=\inf_{a\in A_t(x)}Q_t(h;x,a),
\]
and $V_T^*=g$, $V_t^*=T_tV_{t+1}^*$. The benchmark is the original
fully observed optimum. An acquired-state policy remains feasible for this
benchmark; a restriction to cellwise actions may, however, impose a positive
approximation floor. No vanishing information error is assumed.

Starting from a feasible $\pi^0$, at pass $k$ let $h^k$ have certified
signed policy-evaluation bands
\begin{equation}\label{eq:sweep-bands51}
 a_t^k\leq J_t^{\pi^k}(x)-h_t^k(x)\leq b_t^k,
 \qquad w_t^k=b_t^k-a_t^k.
\end{equation}
These can be obtained from the signed-residual recursion already proved;
the approximator need not have a small absolute level error. Let a feasible
proposal $v_t^k$ satisfy, uniformly in $x$,
\begin{equation}\label{eq:sweep-greedy51}
 Q_t(h_{t+1}^k;x,v_t^k(x))\leq (T_th_{t+1}^k)(x)+\alpha_t^k.
\end{equation}
For its centered comparison with the exact incumbent, define
$d_t^k(x)=Q_t(h_{t+1}^k;x,v_t^k(x))-
Q_t(h_{t+1}^k;x,\pi_t^k(x))$.
A quantitative upper enclosure satisfies
\begin{equation}\label{eq:sweep-enclosure51}
 d_t^k(x)\leq\overline d_t^k(x)\leq d_t^k(x)+\zeta_t^k.
\end{equation}
Here $\zeta_t^k$ includes acquisition-cell variation and integration or
arithmetic enclosure width, not just the machine unit roundoff. The new
policy uses the proposal exactly when
$\overline d_t^k(x)+\beta_tw_{t+1}^k\leq0$ and otherwise uses
$\pi_t^k(x)$ unchanged. For an acquired cell, feasibility and this comparison
must hold throughout that cell, with a single implementable action and
acceptance decision. Bounds \eqref{eq:sweep-greedy51}--\eqref{eq:sweep-enclosure51}
are required on the same cells, including their boundaries.

\begin{theorem}[Finite-sweep accuracy with incumbent preservation]
\label{thm:finite-sweeps51}
Suppose \eqref{eq:sweep-bands51}--\eqref{eq:sweep-enclosure51} hold at each
pass, with nonnegative finite $\alpha_t^k$ and $\zeta_t^k$. Put
\[
 \epsilon_t^k=\alpha_t^k+\zeta_t^k+2\beta_tw_{t+1}^k,
 \qquad E_t^k=\|J_t^{\pi^k}-V_t^*\|_\infty.
\]
Then $J_t^{\pi^{k+1}}\leq J_t^{\pi^k}$ at every date and state, and
\begin{equation}\label{eq:sweep-recursion51}
 E_t^{k+1}\leq\beta_t E_{t+1}^k+\epsilon_t^k,
 \qquad E_T^k=0.
\end{equation}
Set $B_{t,0}=1$ and $B_{t,j}=\prod_{r=t}^{t+j-1}\beta_r$.
For $1\leq m\leq T-t$,
\begin{equation}\label{eq:sweep-unroll51}
 E_t^m\leq B_{t,m}E_{t+m}^0+
 \sum_{j=0}^{m-1}B_{t,j}\epsilon_{t+j}^{m-1-j}.
\end{equation}
In particular, if $\epsilon_t^k\leq\bar\epsilon_t$ uniformly in $k$,
then after $T$ passes
\begin{equation}\label{eq:sweep-final51}
 E_t^T\leq\sum_{j=0}^{T-t-1}B_{t,j}\bar\epsilon_{t+j}.
\end{equation}
With exact policy evaluation up to a datewise constant, exact greedy
proposals, and exact centered comparisons, the policy is optimal after at
most $T$ passes, and no intermediate pass raises its cost.
\end{theorem}

The hypotheses separate three computational tasks. A trained continuation
must be evaluated relative to the policy actually retained; the action
search must approximate the original feasible infimum; and the comparison
must be valid over the available information set. A guarantee for only one
of these tasks does not imply the theorem. The same conclusion applies to
native and neural representations when they implement the same objects.
In particular, it does not turn their common compilation saving into a
neural-specific advantage. Safe policy improvement is an established
principle; the result here supplies the signed-width and acquisition-aware
error account for the incumbent-preserving NBO construction.

\begin{proposition}[An implementable cellwise proposal bound]
\label{prop:cell-greedy51}
Let an acquired cell $C$ have representative $x_C$ and radius $r_C$.
Suppose $Q_t(h;\cdot,a)$ is $L_t$-Lipschitz on the cell for every menu
action, and $Q_t(h;x,\cdot)$ is $D_t$-Lipschitz on feasible actions.
Let $M_C\subseteq\bigcap_{x\in C}A_t(x)$ be a nonempty finite menu such that every
$a\in A_t(x)$, for every $x\in C$, lies within $\delta_{A,C}$ of some
member of $M_C$. If $v_C\in M_C$ is an $\eta_C$-minimizer at $x_C$,
then it satisfies \eqref{eq:sweep-greedy51} on $C$ with
\[
 \alpha_t^k\leq \eta_C+2L_tr_C+D_t\delta_{A,C}.
\]
\end{proposition}
The directed-cover assumption is substantive when constraints move with the
state. It is not implied by a dense grid at the representative alone. A
failed cover or a cell that cannot certify a comparison retains its incumbent;
it must not be counted as a successful accuracy refinement.

\begin{proposition}[Sharpness of the evaluation-width coefficient]
\label{prop:sharp-width51}
Under the theorem's assumptions, the coefficient two on
$\beta_tw_{t+1}^k$ in \eqref{eq:sweep-recursion51} cannot be uniformly
replaced by a smaller constant, even for a one-period deterministic problem
with two actions, an exact greedy proposal, and exact contrast arithmetic.
\end{proposition}

The theorem is constructive but conditional on verified error budgets. It
does not assert that a continuous neural fit uniformly approximates a
discontinuous acquired-policy value, or that refining a fixed acquisition
contract eliminates its information floor. The executed investment study
below certifies a final-date improvement. It is not an execution of $T$
full improvement passes, nor numerical evidence of the finite-sweep
near-optimality conclusion. Exact finite-state regression tests exercise the
new identities, rejection cases, and sharpness example separately from the
economic benchmark.

\section{Allocating resources to an economic tolerance}\label{sec:resources49}
A single resolution parameter imposes an economic restriction that the Bellman problem does not impose. State coverage, action search and integration have different error coefficients and different effects on work. Refining all three at the same rate can therefore spend most of the service budget on the wrong source of error. We separate these choices while retaining the continuation, feasible-witness interpretation and all-state comparison developed above.

\subsection{A separated primitive budget}
Consider the investment family specified in Section~\ref{sec:study49}, with $d\in\{2,3,4\}$ state coordinates, discount $\beta=15/16$, and uniform coordinate grids with $N$ intervals. At a state node use $K+1$ feasible action fractions. Integrate the common scalar innovation using $M$ equal-probability bins and their midpoints. The three integers need not agree. Let $L_t$ be the primitive graph modulus of the constructed witness continuation, and let $D_t$ be its action modulus. A downward numerical displacement of the action grid is denoted by $\delta_a$.

The state and action covering radii satisfy $h_t=d/(2N)$ and $k_t\le 1/(8K)+\delta_a$. The innovation lies in an interval of length $1/16$. Within a bin its expected absolute distance from the midpoint is $1/(64M)$. Since changing the scalar innovation by $z-z'$ changes the next state by at most $d|z-z'|$ in the $\ell^1$ norm, a certified numerical query has an error allowance
\begin{equation}\label{eq:query49}
 e_t\le\frac{\beta d L_{t+1}}{64M}+\xi_t.
\end{equation}
Here $\xi_t$ covers arithmetic, interval evaluation and label rounding. It does not suppress the integration error. The terminal label allowance is $e_T$.

\begin{proposition}[Separated-resource policy account]\label{prop:budget49}
Under the feasible-witness hypotheses, the returned policy satisfies
\begin{equation}\label{eq:separated49}
 0\le J^\pi-V^*\le \frac{a}{N}+\frac{b}{K}+\frac{q}{M}+\eta,
\end{equation}
where
\begin{align}\label{eq:coefficients49}
 a&=d\sum_{t<T}\beta^tL_t+d\beta^TL_g,&
 b&=\frac18\sum_{t<T}\beta^tD_t,\\
 q&=\frac d{32}\sum_{t<T}\beta^{t+1}L_{t+1},&
 \eta&=\sum_{t<T}\beta^t(2\xi_t+D_t\delta_a)+2\beta^Te_T.\notag
\end{align}
The primitive coefficients are independent of the computed labels. The leading number of innovation evaluations in target formation is
\begin{equation}\label{eq:workseparated49}
 T(N+1)^d(K+1)M.
\end{equation}
The compiled continuation additionally incurs its corner-evaluation and preprocessing costs. Pilot selection, policy extraction, verification, arithmetic bit lengths and output are separate charged components.
\end{proposition}
\begin{proof}
Insert the state and action radii and \eqref{eq:query49} into the feasible-witness loss bound. The selected-policy upper account contributes the second copy of $e_t$, not a second action-search allowance. Collecting the coefficients of $N^{-1},K^{-1},M^{-1}$ gives \eqref{eq:coefficients49}. There are $(N+1)^d$ state nodes, $K+1$ actions per node and $M$ innovation representatives per action at each nonterminal date. No independence of fitting errors is used.
\end{proof}

The account explains why a common further state rung need not be a diagonal refinement. In the $T=3,p=4,d=2$ cell, the exact primitive coefficients are approximately $a=51.1569$, $b=2.3650$ and $q=1.0486$. The choice $(N,K,M)=(128,16,8)$ has account about $0.6786+\eta$, below target one when the numerical remainder fits the remaining budget. The diagonal choice $(128,128,128)$ is more accurate, but its target-generation innovation count at that same state grid is $129\cdot128/(17\cdot8)$ times larger. This is an arithmetic comparison of two declared allocations, not an observation of the unexecuted diagonal service or a runtime ratio. Both witness and conventional FVI receive the separated allocation in the experiment.

\subsection{The optimal continuous allocation and its boundary cases}
A useful planning problem minimizes the leading product $N^dKM$ subject to \eqref{eq:separated49}. This proxy isolates the primitive evaluation term; it is not the entire operation count or the measured wall clock. Write $E=\varepsilon-\eta>0$ and suppose first that $a,b,q$ are positive and all optimal resources exceed their required lower bounds.

\begin{proposition}[Error shares and clipped allocation]\label{prop:allocation49}
The unique interior minimizer of the continuous proxy has
\begin{equation}\label{eq:shares49}
 \frac aN=\frac d{d+2}E,\qquad
 \frac bK=\frac1{d+2}E,\qquad
 \frac qM=\frac1{d+2}E.
\end{equation}
For lower resource bounds $r_i\ge\underline r_i>0$, put
$c=(a,b,q)$, $p=(d,1,1)$ and $\bar x_i=c_i/\underline r_i$.
If $E<\sum_i\bar x_i$, there is a unique positive $\lambda$ such that
\begin{equation}\label{eq:clipped49}
 x_i=\min\{\bar x_i,p_i\lambda\},\qquad \sum_i x_i=E;
 \quad r_i=c_i/x_i.
\end{equation}
These resources minimize the constrained continuous proxy. If $E\ge\sum_i\bar x_i$, all lower resource bounds suffice. Coordinates with zero error coefficient remain at their resource lower bound. When dyadic lower bounds are used, rounding each optimal resource upward to a power of two preserves the error guarantee and multiplies this proxy by less than $2^{d+2}$.
\end{proposition}
\begin{proof}
Set $x_i=c_i/r_i$. Minimizing the resource product is equivalent to maximizing $\sum_i p_i\log x_i$ over $0<x_i\le\bar x_i$ and $\sum_i x_i\le E$. Strict concavity and the multiplier conditions give $p_i/x_i$ equal on unconstrained coordinates and larger on capped coordinates. This is precisely \eqref{eq:clipped49}. The sum of its right sides is strictly increasing until all caps bind. The unconstrained case yields \eqref{eq:shares49}. Upward dyadic rounding never enlarges an error term and multiplies $N^dKM$ by less than the stated factor. Neither argument bounds rational-arithmetic bit cost or computer time.
\end{proof}

The optimization is a standard separable resource-allocation calculation. Its role here is to expose the error coefficients of the actual own-future policy construction, including integration and implementation, rather than to claim novelty for the multiplier argument. A finite experimental menu need not contain the continuous optimum. The prototype uses a common, frozen dyadic menu for all generators and records all of it. It does not describe this menu as a globally optimal tuning policy for either FVI or the witness method.

\subsection{A complete tolerance frontier}
For method $m$, let a prospectively fixed ladder return certificates $g_{m,r}$ and cumulative costs $w_{m,r}$, including failed earlier rungs. Define
\begin{equation}\label{eq:frontier49}
 r_m(\varepsilon)=\min\{r:g_{m,r}\le\varepsilon\},\qquad
 W_m(\varepsilon)=w_{m,r_m(\varepsilon)}.
\end{equation}
An empty set is a capped failure, not an interpolated success. Sort the finite union of all certificate values and zero. On each left-closed, right-open interval between successive values, every method's first crossing and complete work are constant. Evaluating the first-crossing rule at the left endpoint therefore reconstructs the entire positive-tolerance frontier exactly, even when successive certificates are not monotone. An unbounded final interval is retained. Repetitions have separate costs, while their mathematical checkpoints must be checked for identity.

This finite partition supplies the appropriate interpretation of a threshold advantage. At a fixed resolution the witness method may cost more. At a tolerance where it crosses one rung earlier, the smaller certificate may more than offset that cost. Whether this happens, and over which tolerance interval, is determined by \eqref{eq:frontier49}; it is not determined by a favorable single threshold. Under an additional homogeneous model $g_m(n)=A_m/n$ and $w_m(n)=C_m n^\kappa$, the continuous frontier ratio is $(C_W/C_F)(A_W/A_F)^\kappa$. This last expression is a conditional explanation, not a fitted scaling law for the experiments. The release supplies the exact finite partition, sensitivity around one, complete prefixes and separate fixed-resolution observations.

\section{From returned policies to economic resource decisions}\label{sec:decisions49}
An all-state upper bound and a comparison of actual policy costs answer different economic questions. We now connect the policies returned by a construction frontier to direct cost inference and then to a choice of implementation resources. The construction data select the policies; independently generated evaluation paths assess them. This separation permits a valid comparison of first-crossing policies with different resolutions.

\subsection{Simultaneous inference after construction}
Let $\mathcal H$ denote all construction information: primitive inputs, numerical labels, fitted continuations, returned policies and complete certificate ladders. Conditional on $\mathcal H$, fix a finite catalogue of implemented-policy contrasts. Each implementation includes its observation rule, original witness or nearest-node selector, true-state feasible repair and action quantization. For a contrast $i$, let $D_i$ be the difference in expected discounted costs under a specified initial law. A common-innovation path score evaluates the two policies on their own endogenous trajectories. It is not the difference between two regret bounds.

Suppose the initial and innovation uniforms are partitioned into equal bins. Independent uniform bin indices are drawn, and interval simulation encloses the score for every continuous realization inside the drawn bins. All unresolved state, selector and repair branches remain in the enclosure. If $[A_i,B_i]$ is a deterministic support interval, clip the path endpoints to that support. Write $\underline m_i,\overline m_i$ for outward sample-mean bounds on the lower and upper endpoint samples and $v_{i,L},v_{i,U}$ for upper bounds on their unbiased sample variances. For $n\ge2$ define
\begin{equation}\label{eq:radius49}
 R_i(v)=\sqrt{2v\ell/n}+\frac{7(B_i-A_i)\ell}{3(n-1)},
 \qquad \ell\ge\log(4K/\alpha),
\end{equation}
where $K$ bounds the number of sampled contrasts.

\begin{theorem}[Construction-conditional economic intervals]\label{thm:conditional49}
Under the stated independent-bin model, conditional on any valid construction history $\mathcal H$, with probability at least $1-\alpha$ all catalogue contrasts satisfy
\begin{equation}\label{eq:direct49}
 D_i\in[l_i,u_i]:=
 [\max\{A_i,\underline m_i-R_i(v_{i,L})\},
  \min\{B_i,\overline m_i+R_i(v_{i,U})\}].
\end{equation}
The same coverage holds unconditionally. Subsequent selection among these covered contrasts, or evaluation at any deterministic implementation-price vector, needs no additional union bound. This conclusion applies only to policies and initial laws included in the catalogue; it does not create direct-cost observations for every unstudied tolerance or initial state.
\end{theorem}
\begin{proof}
For each drawn bin sequence, conditional integration over its continuous realizations places the exact path expectation between the interval endpoints. Uniform mixing over the equal bins recovers the original continuous law. The lower endpoint expectation is thus no larger than $D_i$, and the upper endpoint expectation is no smaller. Apply the bounded empirical-Bernstein inequality of \citet{maurer2009} to the endpoint samples and take the union bound, with outward arithmetic included in the mean, variance and radius. Conditional validity for every $\mathcal H$ implies unconditional validity by integration. On this one simultaneous event, all ensuing deterministic comparisons of the covered quantities hold together.
\end{proof}

In particular, selecting the first policy to meet a specified certificate is permitted when that selection uses construction records rather than the subsequent evaluation paths. Repeated construction clocks are not independent draws from a neural-optimizer population. The statistical randomness here belongs to the evaluation design. A fixed pseudorandom stream implements a reproducible experiment; it is not itself a proof of the independent-bin assumption.

\subsection{Replacement-fee identification}
Write $D=J^W-J^F$. If witness is installed, replacing it by FVI changes total cost by $k-D$, where $k\ge0$ is the replacement fee. If FVI is installed, the change is $k+D$. Thus the two directional gains before fees are $D$ and $-D$.

\begin{corollary}[The break-even fee frontier]\label{cor:fee49}
On an event where $D\in[l,u]$, the two directional gains belong to $[l,u]$ and $[-u,-l]$. The unknown absolute cost difference satisfies
\begin{equation}\label{eq:absolute49}
 |D|\in[\operatorname{dist}(0,[l,u]),\max\{|l|,|u|\}].
\end{equation}
For every $k\ge\max\{0,u,-l\}$ neither replacement strictly lowers expected cost. By contrast, $k\ge\max\{0,l,-u\}$ only removes a \emph{certified} strictly profitable replacement. The two thresholds generally differ. Strict inequalities in the first criterion certify that both replacements strictly increase cost.
\end{corollary}
\begin{proof}
Subtract the gain intervals from $k$ and compare the resulting interval endpoints with zero. Taking the image of $[l,u]$ under the absolute-value map gives \eqref{eq:absolute49}. The assertions hold simultaneously for every nonnegative fee on the original coverage event.
\end{proof}

The fee frontier is more informative than a conclusion at one generous charge. The earlier $1/64$ replacement decision is retained with its original interpretation and confidence account. The new analysis reports the two directional gain intervals and both thresholds instead of selecting a favorable fee after seeing the data. There is no welfare calibration hidden in the units: a decision maker must supply the economically relevant charge, or use the full frontier to see what remains unidentified.

\subsection{Observation costs on the same bit grid}
Let $\mathcal B=\{4,\ldots,16\}$ be the available bit precisions. For one fixed returned controller, write $J_b$ for the expected cost of its actual $b$-bit, robustly repaired and quantized implementation. Use $b_0=16$ as a common reference and obtain simultaneous intervals for $D_b=J_b-J_{b_0}$. The identity interval at $b=b_0$ is exactly $[0,0]$. Let $s_b=d b\sum_{t<T}B_t$ be the discounted number of coordinate-bits purchased; the net cost is $J_b+\lambda s_b$.

\begin{theorem}[Certified choice of actual net implementation cost]\label{thm:net49}
Suppose all $D_b\in[l_b,u_b]$ simultaneously. For any $\lambda\ge0$ define
\begin{align}\label{eq:net49}
 L(\lambda)&=\min_b\{l_b+\lambda s_b\},&
 U(\lambda)&=\min_b\{u_b+\lambda s_b\},\\
 \widehat b(\lambda)&=\min\arg\min_b\{u_b+\lambda s_b\}.\notag
\end{align}
Then, simultaneously for all $\lambda\ge0$,
\begin{equation}\label{eq:netregret49}
 0\le J_{\widehat b}+\lambda s_{\widehat b}
       -\min_{b\in\mathcal B}(J_b+\lambda s_b)
 \le U(\lambda)-L(\lambda).
\end{equation}
The lower and upper functions are finite piecewise-affine envelopes. All changes in their active lines occur at nonnegative intersections of catalogue lines, with exact comparisons deciding ties. The smallest selected precision $\widehat b(\lambda)$ is weakly decreasing in $\lambda$.
\end{theorem}
\begin{proof}
Subtract the common $J_{b_0}$. The selected implementation has net difference at most $U(\lambda)$, and the best catalogue implementation has net difference at least $L(\lambda)$. This proves the regret bound pointwise on the one simultaneous event and hence for every $\lambda$. A minimum of finitely many affine functions changes its active affine expression only at an intersection. For $\lambda_2>\lambda_1$, sum the two minimizing inequalities to obtain $(\lambda_2-\lambda_1)(s_{\widehat b(\lambda_2)}-s_{\widehat b(\lambda_1)})\le0$. Since $s_b$ increases strictly with $b$, the precision comparative static follows.
\end{proof}

This result makes a statement about \emph{actual} net cost within the declared bit catalogue. It differs from minimizing a uniform all-state loss bound. Both criteria are useful: the latter gives distribution-free protection relative to the full-information optimum, while \eqref{eq:netregret49} uses a specified initial law to compare feasible implementations of one returned policy. Neither criterion gives an optimum over arbitrary sensor technologies or all policies.

For a fair comparison the conventional actor receives its own acquisition allowance. Suppose its nearest-node policy has a certified selected-policy account at a reported state, obtained from the graph modulus $L_t^{\mathrm{gr}}$, and its actual interpolated critic has modulus $L_t^f$. Transporting the same selected node action to the true state, then using cell-wide capacity repair, adds at most
\begin{equation}\label{eq:fvisensor49}
 (L_t^{\mathrm{gr}}+L_t^f+2D_tK_t^A)r(b)+D_t\Delta_t.
\end{equation}
Indeed, transport the feasible node pair directly to the true state, bound its node distance by the reported-state distance plus $r(b)$, and bound $f_t(\widehat x)-f_t(x)$ by $L_t^fr(b)$. Robust repair adds $D_t(2K_t^Ar(b)+\Delta_t)$. This proves \eqref{eq:fvisensor49} without assuming actor continuity. For a witness critic $L_t^f=L_t^{\mathrm{gr}}$, it reduces to the preceding witness allowance. The empirical comparison therefore evaluates the all-state-account minimizer, the interval upper-net-cost selector, and a descriptive midpoint selector on the \emph{same} integer bit and price grid. The midpoint selector is not called the true optimum.

\section{Joint policy-cost and resource decisions}\label{sec:joint50}
Let a finite catalogue $\mathcal P$ contain implemented policies, including incumbents, repaired policies and conventional alternatives. Construction, verification and direct evaluation are distinct charged activities. Denote their measured deployment-relevant resource vector by $r_j$; the published timing convention must specify whether an evaluation shared by several candidates is a common fixed charge or an attributable marginal charge. In particular, a joint paired evaluation cannot be counted as zero or silently charged only to the preferred generator.

Conditional on construction information, suppose a simultaneous event satisfies
$J_j\in[L_j,U_j]$ for every $j\in\mathcal P$. For nonnegative prices $\theta$, let $k_j(\theta)$ be a known implementation charge, for example $\theta\cdot r_j$ plus an information fee. Define
\begin{equation}\label{eq:jointselector50}
 \widehat j(\theta)\in\arg\min_{j\in\mathcal P}
   \{U_j+k_j(\theta)\},\qquad
 R(\theta)=U_{\widehat j}+k_{\widehat j}(\theta)
            -\min_{j\in\mathcal P}\{L_j+k_j(\theta)\}.
\end{equation}
A fixed lexicographic tie rule makes the selection measurable.
\begin{proposition}[Actual net-cost frontier]\label{prop:joint50}
On the simultaneous event, for every nonnegative price vector, including one selected after observing the intervals,
\[
 0\le J_{\widehat j}+k_{\widehat j}(\theta)
       -\min_j\{J_j+k_j(\theta)\}\le R(\theta).
\]
A proposed replacement of incumbent $i$ by $j$ is certified nonworsening in net cost whenever a simultaneous upper bound on $J_j-J_i$, plus $k_j(\theta)-k_i(\theta)$ and the replacement fee, is nonpositive. Otherwise a rule retaining $i$ remains safe. Policy-specific certificate widths can be displayed alongside this frontier, but do not determine its ordinate.
\end{proposition}
\begin{proof}
Bound the selected true cost above by its upper line and the best catalogue cost below by the minimum lower line. All inequalities hold on the same finite event for every price, so a new price query does not consume another probability allocation. The replacement assertion applies the paired upper bound directly.\end{proof}

Absolute intervals can be considerably wider than centered paired differences. For a decision between an incumbent and its final repair, the centered interval is therefore the appropriate replacement input. A statistically unresolved interval is not evidence of equivalence. In contrast, Corollary~\ref{cor:terminal50} supplies a deterministic nonincrease, although it does not necessarily certify a strictly positive improvement large enough to cover an implementation fee.

\subsection{Information technology and the price of conservatism}\label{sec:technology50}
An explicit theoretical observation technology gives the bit price units. Suppose a device reports a dyadic cell in each coordinate at date $t$. With $b$ bits per coordinate it transmits $db$ bits and requires $\tau db$ resource units, where $\tau$ is a technological conversion factor. At unit resource price $w$, and with a setup cost $K_b$, its present cost is
\begin{equation}\label{eq:technology50}
 C_b(w,\tau)=K_b+w\tau db\sum_{t<T}B_t.
\end{equation}
Each reported cell entails its own feasible action contract. In the retained sensor experiment, $K_b=0$ and the normalized price $\lambda=w\tau$ takes the declared values. This supplies a structural interpretation and a sensitivity variable, not an estimate of actual device costs. The dataset contains no empirical measurements from which $w$ or $\tau$ could be calibrated.

Theorem~\ref{thm:net49} and Proposition~\ref{prop:joint50} apply to the finite bit catalogue with charges \eqref{eq:technology50}. Their upper-minus-lower envelope is a valid bound on the cost of conservative choice relative to the unknown true catalogue minimizer. Replacing that minimizer by the lowest observed sample mean and calling it the ex post true optimum is not justified. The retained exact price envelopes and their regret bounds answer the identified decision question without making this substitution. A richer technology with fixed charges, nonlinear bandwidth costs or precision-dependent delays can be assessed by changing the known charge functions, but new unobserved policies require new cost bounds.

A similarly explicit interpretation applies to terminal repair. A controller that incurs cost $K$ to install the additional gate, and $q$ per deployed final action, is worth adopting only when its expected discounted policy gain covers those charges. The published centered gain interval defines the entire break-even range. Its mathematical safety does not by itself imply economic adoption at positive computation prices.

\section{Certified computation and economic decisions}\label{sec:study49}
\subsection{The controlled investment family}
There are $d\in\{2,3,4\}$ capital coordinates on $[0,1]^d$, one investment action, and a common continuous shock. The feasible set is
\[
 0\le a\le\kappa(x)=\frac18+\frac1{8d}\sum_{j=1}^d x_j.
\]
With cyclic coordinate indices, $\gamma_j=1/2$ for odd $j$ and $1/4$ for even $j$, and $z\sim\mathrm{Unif}[-1/32,1/32]$, the dynamics are
\begin{equation}\label{eq:investment49}
 F_j(x,a,z)=\frac1{16}+\frac{x_j}2+\frac{x_{j+1}}8
       +\frac{x_j(1-x_{j+1})}{16}+\gamma_ja+(-1)^{j-1}z.
\end{equation}
Current costs are
\begin{equation}\label{eq:cost49}
 c(x,a)=\frac2d\sum_j(x_j-5/8)^2+
 \frac1{4d}\sum_j(x_j-x_{j+1})^2+
 2\left(\frac12-\frac2d\sum_jx_j\right)_+^2+pa^2+4a^4.
\end{equation}
Terminal cost doubles the first coefficient and retains the imbalance and shortage terms, with no action cost. Capital enters both adjustment incentives and feasible capacity. The nonlinear transition and shortage penalty couple the coordinates; the model is not separable linear-quadratic control. Prices $p\in\{1,4\}$ and horizons $T\in\{2,3\}$ are theoretical primitives, not calibrated estimates.

For the declared dimensions, valid primitive moduli in $\ell^1$ are
\[
 C_x=\frac{13}{2d},\quad L_g=\frac9d,\quad F_x=\frac{11}{16},\quad
 F_a=\sum_j\gamma_j,\quad C_a=\frac p2+\frac14,\quad K^A=\frac1{8d}.
\]
The transition is invariant by direct interval bounds. The state column sums of its Jacobian are at most $11/16$. The costs have piecewise-affine gradients; the cube vertices and the shortage-boundary vertices give the stated bounds. The supplemental proof restricts this explicit formula to the declared dimensions rather than presuming it for an arbitrary number of coordinates. The general constructive theorem uses whichever valid primitive moduli the economy supplies.

\subsection{Matched services and realized adaptation}
The execution-before protocol fixes three generators: compiled witness, uniform multilinear FVI, and curvature-refined coordinate FVI. Every generator forms its own future labels and receives the same state-node count, feasible action fractions and integration bins. For each two-state price--horizon cell the common resource ladder is
\[
 (N,K,M)=(4,2,1),(8,2,1),(16,4,2),(32,4,2),(64,8,4),(128,16,8).
\]
The entire ladder is run in three isolated processes per cell and method, rotating method order. This yields 36 two-state services and 216 rungs. State dimension three uses four declared rungs through $(32,8,4)$; dimension four uses three through $(16,8,4)$, with the witness and uniform generators, both horizons, price one and three repetitions. Those additional 24 services contain 84 rungs. The comparison across dimensions uses the same target five, not a dimension-specific certificate threshold. This is a finite common-accuracy experiment at a coarse tolerance, not evidence of an asymptotic or dimension-free rate.

The curvature comparator obtains an own-future pilot on a nine-node coordinate grid. For each axis it takes maxima of adjacent exact second differences across the other coordinates. It then bisects the interval maximizing squared length times the largest pilot curvature over that interval, with deterministic leftmost ties. A negligible positive floor permits a defined rule when the pilot is affine; the dominant unit baseline of the earlier heuristic is absent. Pilots, allocation and the realized axes are stored and charged. Realized nonuniformity, rather than the method's name, determines whether the experiment supplies an adaptive comparison. This remains a coordinate tensor algorithm; it is not presented as a sparse-grid or general adaptive-cell method.

Each rung starts from primitives. The internal complete-prefix clock includes own-future integration, all earlier failed rungs, pilots, exact compilation, actor construction, all-state certification, deployment checks and durable checkpoints. Separate process clocks include startup and the identical non-economic warm-up. A single CPU is selected, numerical-library threads are fixed to one, and CPU frequency is not controlled. Peak resident memory includes the process. Primitive query counts are component counts, not complete FLOPs or bit-operation counts. Independent policy evaluation has a separate joint service and is never added to an unrelated historical clock.

{\small
\setlength{\tabcolsep}{3.5pt}
\begin{longtable}{@{}lrrrr@{}}
\caption{Block A separated-resource attainment}\label{tab:attainment49}\\
\toprule
Method & $4$ & $2$ & $1$ & $1/2$ \\
\midrule
\endfirsthead
\toprule
Method & $4$ & $2$ & $1$ & $1/2$ \\
\midrule
\endhead
\bottomrule
\endfoot
W & 12/12 & 12/12 & 12/12 & 6/12 \\
F & 12/12 & 12/12 & 12/12 & 3/12 \\
A & 12/12 & 12/12 & 12/12 & 3/12 \\
\end{longtable}
}
\noindent{\footnotesize W: compiled witness; F: uniform FVI; A: pilot-curvature bisection FVI. Twelve services per method are four economies, each repeated three times. They are not twelve independent economic tasks or optimizer draws.\par}

{\small
\setlength{\tabcolsep}{3.5pt}
\begin{longtable}{@{}llrrr@{}}
\caption{Complete positive-tolerance frontiers in two states}\label{tab:frontier49}\\
\toprule
$(T,p)$ & Tolerance interval & W: $N$; sec. & F: $N$; sec. & A: $N$; sec. \\
\midrule
\endfirsthead
\toprule
$(T,p)$ & Tolerance interval & W: $N$; sec. & F: $N$; sec. & A: $N$; sec. \\
\midrule
\endhead
\bottomrule
\endfoot
(2,1) & $[0,0.423176)$ & -- & -- & -- \\
(2,1) & $[0.423176,0.483393)$ & 128; 10.658 & -- & -- \\
(2,1) & $[0.483393,0.846352)$ & 128; 10.658 & 128; 8.458 & 128; 9.817 \\
(2,1) & $[0.846352,0.962257)$ & 64; 1.458 & 128; 8.458 & 128; 9.817 \\
(2,1) & $[0.962257,1.6927)$ & 64; 1.458 & 64; 1.134 & 64; 1.514 \\
(2,1) & $[1.6927,1.90651)$ & 32; 0.290 & 64; 1.134 & 64; 1.514 \\
(2,1) & $[1.90651,2.79584)$ & 32; 0.290 & 32; 0.269 & 32; 0.399 \\
(2,1) & $[2.79584,3.21615)$ & 16; 0.092 & 32; 0.269 & 32; 0.399 \\
(2,1) & $[3.21615,5.59167)$ & 16; 0.092 & 16; 0.119 & 16; 0.176 \\
(2,1) & $[5.59167,6.13235)$ & 8; 0.034 & 16; 0.119 & 16; 0.176 \\
(2,1) & $[6.13235,10.0042)$ & 8; 0.034 & 8; 0.061 & 8; 0.090 \\
(2,1) & $[10.0042,10.6876)$ & 4; 0.014 & 8; 0.061 & 8; 0.090 \\
(2,1) & $[10.6876,\infty)$ & 4; 0.014 & 4; 0.029 & 4; 0.042 \\
(2,4) & $[0,0.45093)$ & -- & -- & -- \\
(2,4) & $[0.45093,0.506837)$ & 128; 10.689 & -- & -- \\
(2,4) & $[0.506837,0.901859)$ & 128; 10.689 & 128; 8.453 & 128; 9.797 \\
(2,4) & $[0.901859,1.00914)$ & 64; 1.452 & 128; 8.453 & 128; 9.797 \\
(2,4) & $[1.00914,1.80372)$ & 64; 1.452 & 64; 1.134 & 64; 1.509 \\
(2,4) & $[1.80372,2.00023)$ & 32; 0.290 & 64; 1.134 & 64; 1.509 \\
(2,4) & $[2.00023,2.92224)$ & 32; 0.290 & 32; 0.268 & 32; 0.397 \\
(2,4) & $[2.92224,3.31506)$ & 16; 0.092 & 32; 0.268 & 32; 0.397 \\
(2,4) & $[3.31506,5.84447)$ & 16; 0.092 & 16; 0.118 & 16; 0.177 \\
(2,4) & $[5.84447,6.33036)$ & 8; 0.034 & 16; 0.118 & 16; 0.177 \\
(2,4) & $[6.33036,10.3185)$ & 8; 0.034 & 8; 0.061 & 8; 0.091 \\
(2,4) & $[10.3185,10.9232)$ & 4; 0.014 & 8; 0.061 & 8; 0.091 \\
(2,4) & $[10.9232,\infty)$ & 4; 0.014 & 4; 0.029 & 4; 0.043 \\
(3,1) & $[0,0.63562)$ & -- & -- & -- \\
(3,1) & $[0.63562,0.690808)$ & 128; 15.635 & -- & -- \\
(3,1) & $[0.690808,1.27124)$ & 128; 15.635 & 128; 12.420 & 128; 14.176 \\
(3,1) & $[1.27124,1.37488)$ & 64; 2.070 & 128; 12.420 & 128; 14.176 \\
(3,1) & $[1.37488,2.54248)$ & 64; 2.070 & 64; 1.651 & 64; 2.148 \\
(3,1) & $[2.54248,2.72357)$ & 32; 0.401 & 64; 1.651 & 64; 2.148 \\
(3,1) & $[2.72357,4.11404)$ & 32; 0.401 & 32; 0.388 & 32; 0.568 \\
(3,1) & $[4.11404,4.53322)$ & 16; 0.126 & 32; 0.388 & 32; 0.568 \\
(3,1) & $[4.53322,8.22807)$ & 16; 0.126 & 16; 0.171 & 16; 0.252 \\
(3,1) & $[8.22807,8.6154)$ & 8; 0.046 & 16; 0.171 & 16; 0.252 \\
(3,1) & $[8.6154,14.5143)$ & 8; 0.046 & 8; 0.088 & 8; 0.130 \\
(3,1) & $[14.5143,14.9762)$ & 4; 0.019 & 8; 0.088 & 8; 0.130 \\
(3,1) & $[14.9762,\infty)$ & 4; 0.019 & 4; 0.041 & 4; 0.061 \\
(3,4) & $[0,0.678554)$ & -- & -- & -- \\
(3,4) & $[0.678554,0.723386)$ & 128; 15.815 & -- & -- \\
(3,4) & $[0.723386,1.35711)$ & 128; 15.815 & 128; 12.435 & 128; 14.257 \\
(3,4) & $[1.35711,1.44001)$ & 64; 2.085 & 128; 12.435 & 128; 14.257 \\
(3,4) & $[1.44001,2.71422)$ & 64; 2.085 & 64; 1.643 & 64; 2.169 \\
(3,4) & $[2.71422,2.85349)$ & 32; 0.402 & 64; 1.643 & 64; 2.169 \\
(3,4) & $[2.85349,4.31287)$ & 32; 0.402 & 32; 0.390 & 32; 0.569 \\
(3,4) & $[4.31287,4.66727)$ & 16; 0.127 & 32; 0.390 & 32; 0.569 \\
(3,4) & $[4.66727,8.62574)$ & 16; 0.127 & 16; 0.170 & 16; 0.253 \\
(3,4) & $[8.62574,8.88407)$ & 8; 0.046 & 16; 0.170 & 16; 0.253 \\
(3,4) & $[8.88407,15.0203)$ & 8; 0.046 & 8; 0.088 & 8; 0.131 \\
(3,4) & $[15.0203,15.306)$ & 4; 0.019 & 8; 0.088 & 8; 0.131 \\
(3,4) & $[15.306,\infty)$ & 4; 0.019 & 4; 0.042 & 4; 0.061 \\
\end{longtable}
}
\noindent{\footnotesize Each time is the median complete prefix over three isolated repetitions. A dash is a capped failure. Adjacent bands with identical returned rungs are merged. Printed cutoffs are rounded; exact rational endpoints in the deposited frontier determine membership, including equality. The action and integration budgets are uniquely fixed by $N$ in the common ladder. Minima, maxima and counts are deposited, not replaced by these medians.\par}

{\small
\setlength{\tabcolsep}{3.5pt}
\begin{longtable}{@{}rrrrrrrr@{}}
\caption{Repeated common-accuracy comparison: target five, price one}\label{tab:dimension49}\\
\toprule
$d$ & $T$ & W pass & $N$ & sec. & F pass & $N$ & sec. \\
\midrule
\endfirsthead
\toprule
$d$ & $T$ & W pass & $N$ & sec. & F pass & $N$ & sec. \\
\midrule
\endhead
\bottomrule
\endfoot
2 & 2 & 3/3 & 16 & 0.092 & 3/3 & 16 & 0.119 \\
2 & 3 & 3/3 & 16 & 0.126 & 3/3 & 16 & 0.171 \\
3 & 2 & 3/3 & 16 & 2.424 & 3/3 & 16 & 2.074 \\
3 & 3 & 3/3 & 16 & 3.508 & 3/3 & 16 & 3.032 \\
4 & 2 & 3/3 & 16 & 58.569 & 3/3 & 16 & 60.809 \\
4 & 3 & 3/3 & 16 & 82.156 & 3/3 & 16 & 90.099 \\
\end{longtable}
}
\noindent{\footnotesize All dimensions face the same coarse target five; caps and allocations remain explicit. Two-state observations are reused from the main catalogue and not counted twice. This finite comparison does not identify an asymptotic scaling law.\par}

The original block A contains 60 services and 300 rungs. Repeated mathematical checkpoints are byte-identical. The pilot-curvature bisection generator realizes nonuniform axes in 0 of 252 date models; its name alone does not establish effective adaptation. At common successful declared target-by-repetition comparisons, witness is faster than uniform FVI in 13 of 39 and slower in 26. Against pilot-curvature bisection FVI, the corresponding counts are 30 faster and 9 slower out of 39. These counts summarize selected targets; Table~\ref{tab:frontier49} and the exact deposited partition, rather than a binary attainment headline, are the primary work comparison.

The new first-crossing direct-cost intervals identify lower witness cost in 0 cases, higher cost in 93, and leave 3 unresolved. For the separately evaluated favorable R48 work pair, the corresponding counts are 0, 0 and 8. A cheaper first crossing and a sharper bound therefore receive an actual economic comparison rather than an inferred ranking. The detailed intervals and initial actions are retained, including every adverse outcome.

Across the 24 declared controller--price decisions, the upper-net-cost selector is certified strictly cheaper than the all-state-bound minimizer in 24. Its cost-regret upper bounds relative to the best evaluated bit range from 0.00238686 to 0.00332242 in the normalized economic units. Where the comparison does not identify a strict improvement it remains unresolved; midpoint choices are descriptive, not true optima. This is a finite expected-net-cost conclusion under the uniform initial law, not a calibrated welfare result or a universal preference for coarse observation.

The direct and sensor study contains 200 sampled contrasts, totaling 52,428,800 interval paths, with a separately recorded joint evaluation time of 923.64 seconds. Identity rows use no simulated paths. This evaluation work is not allocated to one favored generator or hidden inside another method's construction clock.


\subsection{Actual first-crossing policies and replacement fees}
For the four two-state economic cells and each target $4,2,1$, the protocol selects repetition zero's first-crossing witness and uniform FVI policies when both attain. It also selects the particular favorable R48 target-two pair: $T=3,p=4$, witness at $N=32$ and FVI at $N=64$. Those old frozen policies are evaluated as an additional pair; their historical clocks are not reconstructed as part of the new service. Each selected pair is compared under the uniform initial law and point laws at $(1/8,1/8)$, $(1/2,1/2)$ and $(7/8,7/8)$, with full information or eight-bit observations. Finite observations use robust cell-wide capacity repair and downward action quantum $1/4096$.

Every sampled contrast uses 262,144 common interval paths, 40-bit equal initial and innovation bins, and outward mean and variance bounds. All new policy-pair and sensor contrasts share one family of at most 300 with error $1/100$ and log upper bound 13. The independent-bin model supplies the statistical assertion; the fixed PCG64 streams supply reproducibility. Old R48 intervals retain their separate original confidence event. Simultaneous interpretation of both old and new families is at least 98 percent by a union bound, not automatically 99 percent.

{\small
\setlength{\tabcolsep}{3.5pt}
\begin{longtable}{@{}lrrrrrrr@{}}
\caption{Actual costs of the returned policies and fee sensitivity}\label{tab:direct49}\\
\toprule
Pair & $\varepsilon$ & Bits & $n$ & \shortstack{W\\lower} & \shortstack{W\\higher} & Unres. & \shortstack{Sufficient\\fee range} \\
\midrule
\endfirsthead
\toprule
Pair & $\varepsilon$ & Bits & $n$ & \shortstack{W\\lower} & \shortstack{W\\higher} & Unres. & \shortstack{Sufficient\\fee range} \\
\midrule
\endhead
\bottomrule
\endfoot
R49 & 4 & E & 16 & 0 & 14 & 2 & [0.00156,0.1042] \\
R49 & 4 & 8 & 16 & 0 & 15 & 1 & [0.001565,0.1047] \\
R49 & 2 & E & 16 & 0 & 16 & 0 & [0.007903,0.05741] \\
R49 & 2 & 8 & 16 & 0 & 16 & 0 & [0.007853,0.05892] \\
R49 & 1 & E & 16 & 0 & 16 & 0 & [0.003374,0.02665] \\
R49 & 1 & 8 & 16 & 0 & 16 & 0 & [0.003379,0.02688] \\
R48 & 2 & E & 4 & 0 & 0 & 4 & [0.001746,0.001875] \\
R48 & 2 & 8 & 4 & 0 & 0 & 4 & [0.001746,0.001886] \\
\end{longtable}
}
\noindent{\footnotesize E denotes full information. Counts cover every declared initial law and economic cell in the group; they do not rank all initial states. The fee is sufficient to certify that neither directional switch pays. It is not the smaller threshold that merely removes identified profitability. All new intervals share one 99 percent family conditional on the independent-bin model.\par}

The complete numerical supplement reports the directional replacement gains for every pair and the fee sufficient to certify that neither switch pays. It also reports the lower threshold that merely removes an identified profitable switch. The old $1/64$ decision is preserved, but no conclusion depends on treating that fee as calibrated. The deposited exact fee frontiers permit any nonnegative fee to be assessed on the same confidence event. Lower regret certificates are never described as observed policy-cost improvements.

\subsection{Purchasing information on the evaluated grid}
For each first-target-one controller, the protocol evaluates every integer precision $b=4,\ldots,16$ under the uniform initial law. Each policy is compared with its own 16-bit implementation; the reference self-comparison is exactly zero. Prices per coordinate-bit are $1/16384,1/4096,1/1024$. The deterministic all-state-account minimizer and the upper-net-cost selector of Theorem~\ref{thm:net49} use precisely that same grid. Neither is compared with an unobserved continuous optimum.

{\small
\setlength{\tabcolsep}{3.5pt}
\begin{longtable}{@{}llrrr@{}}
\caption{Information choice on the same evaluated bit grid}\label{tab:sensor49}\\
\toprule
$T;p$ & Method & \shortstack{Bound-minimizing\\bits} & \shortstack{Net-selected\\bits} & \shortstack{Net-regret\\bound range} \\
\midrule
\endfirsthead
\toprule
$T;p$ & Method & \shortstack{Bound-minimizing\\bits} & \shortstack{Net-selected\\bits} & \shortstack{Net-regret\\bound range} \\
\midrule
\endhead
\bottomrule
\endfoot
2;1 & W & 16,14,12 & 6,6,4 & [0.002563,0.002596] \\
2;1 & F & 16,14,12 & 7,4,4 & [0.002387,0.002426] \\
2;4 & W & 16,14,12 & 6,6,4 & [0.002691,0.00272] \\
2;4 & F & 16,14,12 & 4,4,4 & [0.002564,0.002564] \\
3;1 & W & 16,14,12 & 7,4,4 & [0.003092,0.003107] \\
3;1 & F & 16,14,12 & 6,4,4 & [0.002992,0.003036] \\
3;4 & W & 16,15,13 & 4,4,4 & [0.003322,0.003322] \\
3;4 & F & 16,14,12 & 4,4,4 & [0.003237,0.003237] \\
\end{longtable}
}
\noindent{\footnotesize Bit triples correspond, in order, to prices $1/16384,1/4096,1/1024$ per coordinate-bit. The net selector minimizes the upper actual-net-cost interval. Its regret bound is against the best bit in the finite catalogue, not an estimated unconstrained optimum. Every price-specific value and all nonnegative-price envelopes are deposited.\par}

For each economic cell, the record additionally gives all nonnegative information prices by exact lower and upper affine envelopes. A retrospective price query within that family is allowed by simultaneous coverage and is not a new experiment. Initial-law dependence remains substantive: a uniform-state expected-cost decision need not be optimal at a particular initial state, and a conservatively optimal all-state bound may rationally purchase more information.

\subsection{The retained evidence and its interpretation}
The new service does not rewrite the original nonlinear training studies. The R42 tighter-target neural attainments, six of twelve scalar and five of six coupled services, remain capped-service outcomes; the R44 residual refinements of seven frozen failures remain diagnostics, not retraining successes. The 24 direct neural-minus-ridge intervals remain unresolved. The R45 constructive comparison still has nine of twelve neural and twelve of twelve spline attainments at its tightest target. Fixed64 and adaptive precision still certify all 36 historical precision services, with adaptive slower in 26 and about 0.98 percent slower in aggregate.

The R48 target-one count, twelve of twelve for witness versus six of twelve for each FVI method, also remains unchanged. At a relaxed tolerance 1.014 the old FVI near misses all attain; the release supplies a complete sensitivity table rather than reclassifying those original failures. The new common $N=128$ state rung uses separately declared action and integration budgets. It is not the unexecuted diagonal $(128,128,128)$ experiment. The original 48 constrained-policy intervals, including six favoring FVI and none favoring witness, remain available with newly exposed fee frontiers. Their adverse point-state pattern motivates inspection of actual policies; it is not removed by a favorable all-state certificate.

The complete development edition, unchanged historical companions and machine-readable preservation map retain all of these results and their proofs. This revision adds a controlled resource allocation and actual implementation decision, not a selectively favorable replacement record.

\section{Cost-directed construction and implementation evidence}\label{sec:study50}
We now execute the distinction between constructing a safe policy and improving its actual economic performance. The new policies use the same nonlinear investment primitives as Section~\ref{sec:study49}; neither the price parameter nor the terminal penalty is changed to favor the proposed action rule. All earlier continuation weights, owners, policies, clocks and adverse contrasts retain their historical identities. The recovered R50 main experiment and common-accuracy amendment had separate local pre-execution freezes. R51 is an explicitly known-design reproduction, with its own Git source identity and runner clocks; the same draws do not provide additional independent observations.

\subsection{Prospective resource selection and a residual-driven comparator}
For the primary block, select powers of two $N\in\{4,8,16,32,64\}$ and $K,M\in\{1,2,4,8,16\}$ before constructing a policy. Among allocations satisfying the primitive witness budget $a/N+b/K+q/M\le\varepsilon-10^{-8}$, minimize the exact integer query proxy $T(N+1)^d(K+1)M$, with lexicographic ties. The reserve is checked against the realized certificate; it is not a universal fixed-precision theorem. This rule supplies an executed version of the separated allocation, rather than a retrospective proxy calculation. Conventional FVI receives the same allocation but its own actual-modulus certificate. The witness planning formula alone does not guarantee FVI attainment.

The catalogue contains seven economic cells: $d=2,3,4$, $T=2,p=1,\varepsilon=5$, and $d=2$, $T=2,3$, $p=1,4,\varepsilon=2$. Witness and uniform FVI are each run on all cells. Two additional services choose the smallest isotropic power-of-two allocation $N=K=M$ satisfying the same primitive target-five budget in $d=2,T=2,p=1$. A further service in that cell uses residual-driven FVI. Each of these seventeen specifications is executed three times in separate single-CPU processes, with identical warm-up and rotated service order. Repetitions estimate local timing dispersion, not a distribution of economic tasks or neural initializations.

The residual-driven comparator queries its own future Bellman objective at interval endpoints and midpoints. It inserts coordinate knots where the largest interpolation surplus across the other coordinates is greatest, then equidistributes square-root surplus mass along each axis, with dyadic rounding and explicit monotone-knot checks. This last step follows the local quadratic interpolation-error scaling, rather than the curvature-mass rule of the earlier graded block. All pilot queries, residual comparisons, compilation and output are charged. Recorded geometry, not the method name, determines whether adaptation occurred. This is a coordinate-tensor comparator, not a claim to have executed the whole adaptive sparse-grid literature.

The original R50 design and source were frozen locally before their measured execution. Its early bisection-only development test and subsequent equidistribution correction are retained with their original records. R51 freezes the recovered algorithm before a known-design reproduction; it is not presented as a new blind experiment. The present clocks belong only to this execution, and the earlier R49 evidence is unchanged.

{\small
\setlength{\tabcolsep}{3pt}
\begin{longtable}{@{}llrrrrrr@{}}
\caption{Primary prospective allocations and realized certificates}\label{tab:services50}\\
\toprule
$d;T;p$ & Method & $N,K,M$ & $\varepsilon$ & Bound & Pass & Sec. & Sec. range \\
\midrule
\endfirsthead
\toprule
$d;T;p$ & Method & $N,K,M$ & $\varepsilon$ & Bound & Pass & Sec. & Sec. range \\
\midrule
\endhead
\bottomrule
\endfoot
2;2;1 & W & 32,2,2 & 2 & 1.9713 & 3/3 & 0.0847 & [0.0841,0.0852] \\
2;2;1 & W I & 8,8,8 & 5 & 4.6296 & 3/3 & 0.0195 & [0.0193,0.0200] \\
2;2;1 & W & 16,1,1 & 5 & 3.9427 & 3/3 & 0.0240 & [0.0238,0.0268] \\
2;2;4 & W & 32,8,1 & 2 & 1.9317 & 3/3 & 0.1086 & [0.1072,0.1105] \\
2;3;1 & W & 64,4,2 & 2 & 1.7567 & 3/3 & 0.5059 & [0.5001,0.5075] \\
2;3;4 & W & 64,4,2 & 2 & 1.9149 & 3/3 & 0.5058 & [0.4962,0.5095] \\
3;2;1 & W & 16,1,1 & 5 & 4.0581 & 3/3 & 0.4424 & [0.4335,0.4610] \\
4;2;1 & W & 8,8,2 & 5 & 4.8628 & 3/3 & 1.8441 & [1.8348,1.9009] \\
2;2;1 & A & 16,1,1 & 5 & 4.2586 & 3/3 & 1.0709 & [1.0636,1.0941] \\
2;2;1 & F & 32,2,2 & 2 & 2.1589 & 0/3 & 0.0644 & [0.0631,0.0646] \\
2;2;1 & F I & 8,8,8 & 5 & 5.3703 & 0/3 & 0.0232 & [0.0227,0.0233] \\
2;2;1 & F & 16,1,1 & 5 & 4.2199 & 3/3 & 0.0243 & [0.0242,0.0246] \\
2;2;4 & F & 32,8,1 & 2 & 2.0923 & 0/3 & 0.0900 & [0.0869,0.0918] \\
2;3;1 & F & 64,4,2 & 2 & 1.7761 & 3/3 & 0.3650 & [0.3495,0.3661] \\
2;3;4 & F & 64,4,2 & 2 & 1.9038 & 3/3 & 0.3606 & [0.3590,0.3723] \\
3;2;1 & F & 16,1,1 & 5 & 4.3316 & 3/3 & 0.3890 & [0.3831,0.4258] \\
4;2;1 & F & 8,8,2 & 5 & 5.653 & 0/3 & 2.4360 & [2.4275,2.4645] \\
\end{longtable}
}
\noindent{\footnotesize W: native witness with equivalent ReLU realization; F: uniform FVI; A: residual-driven FVI; I: isotropic allocation. Each specification has three identical mathematical checkpoints and three distinct clocks. The primitive planning guarantee is witness-specific; FVI uses its actual critic modulus and certificate. Decimal rounding never determines attainment.\par}

The primary block certifies 24 of 24 witness services, 12 of 24 uniform FVI services and 3 of three residual-driven services. These counts include timing repetitions, not independent policy draws. The residual comparator has nonuniform axes in all 3 date models of its repeated checkpoint. Its bound is 4.2586, compared with 4.2199 for uniform FVI at the same quota. Its pilot work is retained.

In the isotropic comparison cell, the planned witness uses 1,156 target innovation evaluations, versus 11,664 for the isotropic service. The latter has fewer state nodes; query savings are therefore not a claim of less work in every category. The table reports measured complete clocks rather than inferring a speed ratio from this count.


\subsection{A separate common-accuracy completion}
Inspection of the primary block showed that sharing a witness-sufficient allocation does not make the FVI certificate pass. We therefore froze a separate, explicit service: in $d=2,3,4$ at $T=2,p=1,\varepsilon=5$, begin with the same planner; if a method's actual certificate fails, double only $N$, reconstruct from primitives and retain the failed attempt. Stop at first success or at $N=64$. Both methods face this rule. There are three new isolated repetitions per method and dimension. The source freeze and results are separate from the primary catalogue, and old failures remain failures.

{\small
\setlength{\tabcolsep}{3pt}
\begin{longtable}{@{}rlrrrrrrr@{}}
\caption{Separate common-target completion: horizon two, price one, target five}\label{tab:common50}\\
\toprule
$d$ & Method & $N$ range & Tries & Bound & Pass & Sec. & Sec. range & MiB \\
\midrule
\endfirsthead
\toprule
$d$ & Method & $N$ range & Tries & Bound & Pass & Sec. & Sec. range & MiB \\
\midrule
\endhead
\bottomrule
\endfoot
2 & W & 16--16 & 1 & 3.9427 & 3/3 & 0.0242 & [0.0223,0.0281] & 41.4 \\
2 & F & 16--16 & 1 & 4.2199 & 3/3 & 0.0244 & [0.0236,0.0275] & 41.7 \\
3 & W & 16--16 & 1 & 4.0581 & 3/3 & 0.4383 & [0.4155,0.4386] & 44.3 \\
3 & F & 16--16 & 1 & 4.3316 & 3/3 & 0.3913 & [0.3840,0.3935] & 42.3 \\
4 & W & 8--8 & 1 & 4.8628 & 3/3 & 1.7418 & [1.7139,1.8030] & 50.5 \\
4 & F & 8--16 & 2 & 3.1181 & 3/3 & 29.0946 & [28.9063,29.3447] & 169.8 \\
\end{longtable}
}
\noindent{\footnotesize All failed state refinements are included in the complete prefix. MiB is the maximum process peak over three repetitions, not just model storage. The amendment followed the primary outcomes and was frozen before its own execution. Neither block overwrites the other. Host exclusivity and CPU frequency were not controlled.\par}

The returned policies receive new direct actual-cost comparisons below. These observations provide a common target in every reported dimension. They remain a finite low-dimensional tensor experiment. The construction and corner factors $(N+1)^d$ and $2^d$ are not removed by the compiler, the terminal repair or this allocation rule.

\subsection{Direct policy costs and the source of removable loss}
The primary cost experiment fixes twenty-five groups. Sixteen use the original first-target-two R49 witness and FVI policies in all four price--horizon cells, evaluated under the uniform law and each of the three declared point laws. Seven use the new planned services under uniform initial states; one compares residual-driven and uniform FVI; one uses the isotropic allocations. A separate three-group family evaluates the common-accuracy amendment's returned policies. Every group contains both original and final-repaired policies, with the same eight-bit observation and $1/4096$ action contract. Both generators receive the same improvement.

Each group uses 131,072 continuous-law interval paths and analytic final-innovation integration. Four absolute costs and all six pair contrasts are recorded. The primary 250 estimands and supplemental 30 estimands are covered by the conservative endpoint union-bound budget of 512 and log upper bound 13. The supplemental groups have their own conditional independent-bin model within the declared simultaneous family; the policies may depend on earlier construction information. Fixed streams establish reproducibility only. The older R49 family retains its own separate error account.

{\small
\setlength{\tabcolsep}{3pt}
\begin{longtable}{@{}lrrr@{}}
\caption{Actual policy differences and removable final-action loss}\label{tab:policy50}\\
\toprule
Case & $J_L-J_R$ & $J_{L^+}-J_{R^+}$ & $J_L-J_{L^+}$ \\
\midrule
\endfirsthead
\toprule
Case & $J_L-J_R$ & $J_{L^+}-J_{R^+}$ & $J_L-J_{L^+}$ \\
\midrule
\endhead
\bottomrule
\endfoot
C2;2;1;U & [0.023033,0.028988] & [0.012051,0.017775] & [0.011448,0.014567] \\
C3;2;1;U & [0.01796,0.023776] & [0.010129,0.015794] & [0.0077535,0.010757] \\
C4;2;1;U & [0.033636,0.039473] & [0.015912,0.021433] & [0.016432,0.019542] \\
I2;2;1;U & [0.042713,0.048868] & [0.018209,0.023856] & [0.023541,0.026849] \\
O2;2;1;M & [0.0064146,0.011415] & [-0.0023843,0.0023843] & [0.007622,0.010239] \\
O2;2;1;L & [0.054972,0.060123] & [0.044007,0.048792] & [0.0099411,0.012722] \\
O2;2;1;H & [0.0042173,0.0090667] & [0.0022567,0.007049] & [0.0014934,0.003946] \\
O2;2;1;U & [0.013161,0.018416] & [0.0064562,0.011605] & [0.0056474,0.0083165] \\
O2;2;4;M & [0.035287,0.040496] & [0.030381,0.035526] & [0.0036933,0.0063836] \\
O2;2;4;L & [0.020345,0.025506] & [0.017976,0.023095] & [0.0010828,0.0036974] \\
O2;2;4;H & [0.039236,0.044497] & [0.034549,0.039692] & [0.0034416,0.006155] \\
O2;2;4;U & [0.011181,0.016658] & [0.0058295,0.011256] & [0.004188,0.0069254] \\
O2;3;1;M & [0.015697,0.021868] & [0.012485,0.018641] & [0.0016956,0.0047777] \\
O2;3;1;L & [0.026367,0.032516] & [0.022928,0.029055] & [0.0020031,0.005116] \\
O2;3;1;H & [0.011221,0.017268] & [0.0095515,0.015586] & [0.00030976,0.0033545] \\
O2;3;1;U & [0.0095127,0.015773] & [0.0065302,0.012761] & [0.0014967,0.0045926] \\
O2;3;4;M & [0.010433,0.016948] & [0.0089337,0.015429] & [9.4671e-05,0.0033779] \\
O2;3;4;L & [0.0095446,0.01608] & [0.0074164,0.013916] & [0.0005048,0.0037993] \\
O2;3;4;H & [0.0073415,0.013898] & [0.0046083,0.011119] & [0.0011154,0.0044219] \\
O2;3;4;U & [0.0056055,0.012273] & [0.0039105,0.010565] & [0.00021625,0.0035063] \\
Q2;2;1;U & [0.01186,0.017187] & [0.0056983,0.010907] & [0.0055893,0.0083122] \\
P2;2;1;U & [0.022852,0.028801] & [0.011835,0.01755] & [0.011497,0.014615] \\
Q2;2;4;U & [0.0086864,0.014074] & [0.004022,0.0093656] & [0.003535,0.0062382] \\
Q2;3;1;U & [0.01136,0.017753] & [0.0081953,0.014553] & [0.0017403,0.0048791] \\
Q2;3;4;U & [0.010559,0.017349] & [0.0074835,0.014248] & [0.0020091,0.0053568] \\
P3;2;1;U & [0.017874,0.023687] & [0.0099988,0.015657] & [0.0078066,0.010814] \\
P4;2;1;U & [0.033196,0.039029] & [0.015676,0.021193] & [0.016374,0.019482] \\
A2;2;1;U & [-0.0024517,0.0023493] & [-0.0024125,0.0023985] & [0.00062241,0.0030972] \\
\end{longtable}
}
\noindent{\footnotesize A case is block/dimension; horizon; price; initial law. O: original R49 target-two policies; P: primary target-five plan; Q: primary target-two plan; C: common-accuracy amendment; I: isotropic; A: surplus FVI versus uniform FVI. Except A, left is witness and right is uniform FVI. U is uniform; L,M,H are point states 1/8,1/2,7/8. Superscript + denotes the identical terminal safety rule. Five-significant-digit display is not an exact endpoint; the rational records govern all conclusions.\par}

All 27 declared witness comparison groups identify a strictly positive reduction of the witness incumbent's actual expected cost. The simultaneous gain lower endpoints range from 9.4671e-05 to 0.023541; upper endpoints range from 0.0033545 to 0.026849. This is a law-specific statement for the stated policies, not strict improvement at every initial state or 27 independent training successes.

Before repair, witness cost is identified as higher in 27 groups. After both methods receive the same repair, it remains higher in 26, is lower in 0, and is unresolved in 1. The residual-driven FVI comparison is unresolved both before and after repair. Thus the terminal construction repairs a genuine economic loss without manufacturing a reversal of the strong conventional comparison.

The active evaluator records 3,670,016 interval paths, 7,340,032 cell-gate queries and 5,447,560 accepted changes across the two incumbents. No acquisition ambiguity or broad-hull fallback is recorded. Its joint direct-evaluation time before durable record output is 127.95 seconds. That total is evaluation work, not a speed comparison between independently executed policies.


For each incumbent the direct reduction is the centered final contrast. The remaining witness--FVI difference after both repairs is an actual comparison of the two changed policies. The identity
\[
 J_W-J_F=(J_W-J_{W^+})+(J_{W^+}-J_{F^+})-(J_F-J_{F^+})
\]
separates removable final-action losses from the difference that remains. It does not attribute every remaining discrepancy to a single representation mechanism. The deposited date-level action means, terminal state means, accepted and blocked proposals, and capacity/quantum repair counts permit this distinction to be checked. A finite-sample positive gain can justify a strictly positive replacement charge up to its simultaneous lower endpoint; a deterministic nonincrease alone does not justify an arbitrary implementation fee.

\subsection{Actual-cost frontiers and complete resource boundaries}
The supplement reports every group's four absolute cost intervals and every direct contrast. For each new construction, its resource coordinate includes the full own-future service and, separately, the joint evaluation charge. The reader can use either a shared fixed evaluation cost for the catalogue decision or the complete joint cost required to certify a particular comparison. We do not claim to have measured independent marginal costs for four separately simulated policies when the implementation shares their earlier trajectories. The pair clock is an explicitly charged object, not an unpriced convenience.

{\small
\setlength{\tabcolsep}{3pt}
\begin{longtable}{@{}llrrrr@{}}
\caption{Construction work and actual cost of new implementations}\label{tab:joint50}\\
\toprule
Case & Method & Build sec. & Original cost & Repaired cost & Pair sec. \\
\midrule
\endfirsthead
\toprule
Case & Method & Build sec. & Original cost & Repaired cost & Pair sec. \\
\midrule
\endhead
\bottomrule
\endfoot
C2;2;1;U & W & 0.0242 & [0.6242,0.6378] & [0.61117,0.62482] & 3.799 \\
C2;2;1;U & F & 0.0244 & [0.59821,0.61177] & [0.59628,0.60988] & 3.799 \\
C3;2;1;U & W & 0.4383 & [0.58992,0.60068] & [0.58062,0.59146] & 6.102 \\
C3;2;1;U & F & 0.3913 & [0.56904,0.57982] & [0.56768,0.57849] & 6.102 \\
C4;2;1;U & W & 1.7418 & [0.62188,0.63199] & [0.60399,0.61391] & 9.790 \\
C4;2;1;U & F & 29.0946 & [0.5855,0.59525] & [0.5854,0.59515] & 9.790 \\
I2;2;1;U & W & 0.0195 & [0.63959,0.65386] & [0.61456,0.6285] & 3.744 \\
I2;2;1;U & F & 0.0232 & [0.59407,0.60781] & [0.59363,0.60737] & 3.744 \\
Q2;2;1;U & W & 0.0847 & [0.60972,0.62332] & [0.60276,0.61638] & 3.870 \\
Q2;2;1;U & F & 0.0644 & [0.59518,0.60882] & [0.59444,0.6081] & 3.870 \\
P2;2;1;U & W & 0.0240 & [0.6222,0.63577] & [0.60912,0.62274] & 3.901 \\
P2;2;1;U & F & 0.0243 & [0.59639,0.60993] & [0.59445,0.60803] & 3.901 \\
Q2;2;4;U & W & 0.1086 & [0.69243,0.70702] & [0.68754,0.70213] & 3.919 \\
Q2;2;4;U & F & 0.0900 & [0.68106,0.69562] & [0.68086,0.69542] & 3.919 \\
Q2;3;1;U & W & 0.5059 & [0.72237,0.73696] & [0.71906,0.73365] & 4.406 \\
Q2;3;1;U & F & 0.3650 & [0.70782,0.7224] & [0.70769,0.72228] & 4.406 \\
Q2;3;4;U & W & 0.5058 & [0.86908,0.88468] & [0.86539,0.881] & 4.513 \\
Q2;3;4;U & F & 0.3606 & [0.85511,0.87074] & [0.85451,0.87015] & 4.513 \\
P3;2;1;U & W & 0.4424 & [0.58981,0.60055] & [0.58046,0.59128] & 6.294 \\
P3;2;1;U & F & 0.3890 & [0.56902,0.57978] & [0.56764,0.57844] & 6.294 \\
P4;2;1;U & W & 1.8441 & [0.61951,0.62957] & [0.60168,0.61155] & 7.713 \\
P4;2;1;U & F & 2.4360 & [0.58357,0.59328] & [0.58332,0.59303] & 7.713 \\
A2;2;1;U & A & 1.0709 & [0.59772,0.61129] & [0.59585,0.60945] & 3.797 \\
A2;2;1;U & F & 0.0243 & [0.59777,0.61134] & [0.59585,0.60946] & 3.797 \\
\end{longtable}
}
\noindent{\footnotesize Pair sec. is a shared cost of the whole four-policy comparison, repeated in the display for clarity, not independently attributable to each row. Construction medians use the same cell and fresh records; no historical clocks are spliced. Absolute-cost intervals are wider than centered pair intervals. Charges enter Proposition~\ref{prop:joint50} with their explicitly chosen shared or marginal accounting convention.\par}

All stored coefficients, original owners, action scalars, model evaluation categories, transform relaxations, interpolation/corner operations, terminal derivative queries, whole-cell tests, exact fallbacks, numerical interval widths, serialized bytes and process peak resident memory remain in the records. The counters are primitive categories, not a complete FLOP or bit-operation count. The latter would require an additional cost model. Direct comparison time is not added to clocks from another host or to a historical failure in order to manufacture a new complete service.

The retained sensor catalogue has a structural interpretation through equation~\eqref{eq:technology50}. Its prices remain normalized primitives, and its confidence-envelope regret is relative to the best observed implementation in the finite catalogue. This revision neither estimates a sensor technology nor substitutes a sample-mean winner for the true expected-net-cost minimizer.

\subsection{Historical findings that the new evidence does not replace}
The earlier 93 of 96 strictly higher witness-cost contrasts remain visible in Table~\ref{tab:direct49}. Original curvature-FVI never changed its grid in the realized block. The later genuinely graded method also remains an adverse conventional-comparator experiment: it was less precise in all twelve matched final services and slower in all twelve. These two findings concern the two specified heuristics, not all adaptive dynamic programming. The present surplus comparator adds an independently recorded realized partition, while uniform FVI remains the strong baseline.

Likewise, a successful final-action repair does not establish better earlier continuation learning, universal neural superiority, a guarantee for unrestricted nonconvex training, or a dimension-free policy construction. What is established is a cost-directed feasible improvement with an explicit safety condition, its executed original-economy specialization, and source-bound resource and actual-cost comparisons. These results extend the NBO policy-evaluation and improvement program without replacing its economic subject or erasing its previous evidence.


\section{Economic scope and conclusion}\label{sec:conclusion49}
The economic contribution is an explicit connection between a policy's construction resources and its implementable cost. A primitive error budget can determine whether to refine the state cover, action search or integration. A complete finite tolerance frontier states when one generator returns a certified policy earlier. Independent comparisons of the returned policies then determine what can be said about controller replacement and purchased information. These are different steps; none can be replaced by the favorable outcome of another.

The constructive backend and exact compiler make positive deterministic statements under explicit moduli and effective-computability conditions. Their guarantees are not success probabilities for unrelated nonconvex training algorithms. The matched experiments and direct-cost intervals make finite, source-bound claims about the declared economies, resource ladders and initial laws. Common-accuracy dimension cells do not establish an asymptotic scaling law, and endogenous precision within a finite sensor catalogue is not an optimum over all information technologies. The original broad theory, applications and adverse experimental record remain part of the complete NBO development.

The cost-directed result adds a further positive statement: a feasible acquired policy can be improved without sacrificing its incumbent cost guarantee, provided its centered advantage is verified. In the executed final-stage specialization every studied witness incumbent has a strictly positive identified reduction under the declared initial law. Most witness--FVI differences remain unfavorable after both methods are repaired. The improvement theorem and this comparative observation are consistent: safe improvement relative to an incumbent is not dominance over another method.

Neural Bellman Operators therefore provide a way to organize continuation construction, feasible implementation, certification and economic resource choice around the policy actually returned. Whether a particular representation or generator is preferable is decided by its full certified frontier and actual economic consequences, rather than by the name of its function class.

\appendix
\section{Submission and preservation map}\label{sec:map49}
This article and its technical supplement are the active R51 exposition of the same paper. Every R49 main-article and supplementary label is retained in its corresponding current document. The complete development editions retain the R48 main and supplementary bodies and add the current theory and evidence, with a reading notice distinguishing historical wording from the active empirical account. The original R49 sources remain unchanged under \texttt{preserved/R49}; R48 sources and earlier theory/application companions are likewise retained. Preservation is not a new mathematical validation of every historical application.

The local revision contains ordinary sources, both immutable R49 evidence archives, the R48 publication input, newly executed records, generated tables, a point-by-point response and compiled documents. The direct numerical correction retains its original observations and replays the same streams with explicitly outward cross-policy subtraction. The corrected intervals, not their raw predecessors, govern the active direct-cost tables. Source and evidence identities are checked before a clean offline Git-archive rebuild. The local delivery record distinguishes a completed build from a remote publication: no remote branch or commit is claimed without a confirmed push. This limitation is administrative, not a substitute for the substantive revised manuscript.

\bibliographystyle{ecta-fullname}
\bibliography{references}
\end{document}
